% Extensions to the bulk-boundary-line Koszul triangle chapter.
% Four sections: M2 full computation, Kodaira-Spencer example,
% global Koszul triangle theorem, Swiss-cheese convolution algebra.

\providecommand{\kk}{\mathbb{k}}
\providecommand{\cO}{\mathcal O}
\providecommand{\SCop}{\mathsf{SC}^{\mathrm{ch,top}}}
\providecommand{\gSC}{\mathfrak{g}^{\mathrm{SC}}}
\providecommand{\cC}{\mathcal C}
\providecommand{\cL}{\mathcal L}
\providecommand{\cM}{\mathcal M}
\providecommand{\cJ}{\mathcal J}
\providecommand{\Ainf}{\mathsf{A}_{\infty}}
\providecommand{\Linf}{L_\infty}
\providecommand{\Sym}{\operatorname{Sym}}
\providecommand{\PV}{\operatorname{PV}}
\providecommand{\ot}{\otimes}
\providecommand{\CE}{\operatorname{CE}}
\providecommand{\HH}{\operatorname{HH}}
\providecommand{\RHom}{\operatorname{RHom}}
\providecommand{\BBar}{\operatorname{Bar}}
\providecommand{\Cobar}{\operatorname{\Omega}}
\providecommand{\Spec}{\operatorname{Spec}}
\providecommand{\Spf}{\operatorname{Spf}}
\providecommand{\End}{\operatorname{End}}
\providecommand{\Hom}{\operatorname{Hom}}
\providecommand{\Der}{\operatorname{Der}}
\providecommand{\gr}{\operatorname{gr}}
\providecommand{\im}{\operatorname{im}}
\providecommand{\coker}{\operatorname{coker}}
\providecommand{\id}{\operatorname{id}}
\providecommand{\Perf}{\operatorname{Perf}}
\providecommand{\Cl}{\operatorname{Cl}}
\providecommand{\dCrit}{\operatorname{dCrit}}
\providecommand{\Zder}{\operatorname{Z}_{\mathrm{der}}}
\providecommand{\mc}{\operatorname{MC}}
\providecommand{\op}{\mathrm{op}}
\providecommand{\der}{\mathrm{der}}
\providecommand{\Abulk}{A_{\mathrm{bulk}}}
\providecommand{\Bbound}{B_{\partial}}
\providecommand{\Kline}{K_{\mathrm{line}}}
\providecommand{\intr}{\mathrm{intr}}
\providecommand{\free}{\mathrm{free}}
\providecommand{\eff}{\mathrm{eff}}
\providecommand{\st}{\mathrm{st}}
\providecommand{\ex}{\mathrm{ex}}
\providecommand{\minmod}{\mathrm{min}}
\providecommand{\expmod}{\mathrm{exp}}
\providecommand{\hkR}{\mathrm{HKR}}
\providecommand{\sn}{\mathrm{SN}}
\providecommand{\lines}{\mathsf{Lines}}
\providecommand{\hSym}{\widehat{\Sym}}
\providecommand{\ModEnv}{\operatorname{ModEnv}}
\providecommand{\Tr}{\operatorname{Tr}}
\providecommand{\ohat}{\widehat{\otimes}}
\providecommand{\Mbar}{\overline{\mathcal M}}
\providecommand{\Rring}{R^{\bullet}}
\providecommand{\ov}[1]{\overline{#1}}
\providecommand{\M}{\mathcal{M}}
\providecommand{\barB}{\overline{B}}
\providecommand{\Diff}{\operatorname{Diff}}
\providecommand{\ad}{\operatorname{ad}}
\providecommand{\tr}{\operatorname{tr}}


%================================================================
% THE M2-BRANE COMPUTATION IN FULL
%================================================================

\section{The \texorpdfstring{$M2$}{M2}-brane computation in full}
% label removed: sec:m2-full-computation

The definitions in \S\ref{subsec:m2-ambient-construction} established the
large-$K$ boundary algebra $A_{M2,\infty}$, its Koszul dual, and the
master lift $\Theta_{M2}$. This section performs the full
computation, making the constructions of the preceding section explicit. The treatment is complete through bar degree~$4$ and
produces explicit formulas for every component of the holographic
modular Koszul datum.

\subsection{The large-\texorpdfstring{$K$}{K} boundary algebra:
generators, filtration, and bracket}
% label removed: subsec:m2-generators-filtration

The ambient algebra is the completed universal
enveloping algebra of the Lie algebra
$\mathfrak{gl}_K \otimes \Diff(\C)$, where $\Diff(\C)$ denotes the
associative algebra of polynomial differential operators on $\C$.
The generators are
\begin{equation}% label removed: eq:m2-generators
E^{(m,n)}_{\alpha\beta}
:= E_{\alpha\beta}\,z^m\partial^n,
\qquad
1\le \alpha,\beta\le K,\quad m,n\ge 0,
\end{equation}
where $E_{\alpha\beta}$ is the elementary matrix with a~$1$ in the
$(\alpha,\beta)$-entry. The weight of a generator is
\begin{equation}% label removed: eq:m2-weight
\operatorname{wt}(E^{(m,n)}_{\alpha\beta}) = m + n,
\end{equation}
and the algebra carries the weight filtration
$F^{\le w} A_{M2,\infty}
= \overline{\bigoplus}_{m+n \le w} \kk\cdot E^{(m,n)}_{\alpha\beta}$.

The completed algebra $A_{M2,\infty}$ is the completion of the universal
enveloping algebra with respect to this filtration. At the
classical level the bracket is the commutator in the associative
algebra $\mathfrak{gl}_K \otimes \Diff(\C)$:

\begin{proposition}[Classical bracket; \ClaimStatusProvedHere]
% label removed: prop:m2-classical-bracket
The classical Lie bracket of two generators is
\begin{equation}% label removed: eq:m2-classical-bracket
[E^{(m,n)}_{\alpha\beta},\, E^{(p,q)}_{\gamma\delta}]_0
\;=\;
\delta_{\beta\gamma}
\sum_{j=0}^{n}\binom{n}{j}
\frac{p!}{(p-j)!}\,
E^{(m+p-j,\,n+q-j)}_{\alpha\delta}
\;-\;
\delta_{\alpha\delta}
\sum_{j=0}^{q}\binom{q}{j}
\frac{m!}{(m-j)!}\,
E^{(m+p-j,\,n+q-j)}_{\gamma\beta},
\end{equation}
where the factorials are interpreted as zero when $j>p$ or $j>m$.
\end{proposition}

\begin{proof}
The product in $\Diff(\C)$ is
\[
(z^m\partial^n)(z^p\partial^q)
= \sum_{j=0}^{n}\binom{n}{j}
\frac{p!}{(p-j)!}\,z^{m+p-j}\partial^{n+q-j}.
\]
The matrix part multiplies by the rule
$E_{\alpha\beta}E_{\gamma\delta}
= \delta_{\beta\gamma}E_{\alpha\delta}$.
Taking the commutator produces the two terms.
\end{proof}

\begin{corollary}[Low-weight brackets; \ClaimStatusProvedHere]
% label removed: cor:m2-low-weight-brackets
At the lowest weights:
\begin{align}
[E^{(1,0)}_{\alpha\beta},\,E^{(0,1)}_{\gamma\delta}]_0
&= -\delta_{\beta\gamma}\,E^{(0,0)}_{\alpha\delta}
+ \delta_{\alpha\delta}\,E^{(0,0)}_{\gamma\beta},
% label removed: eq:m2-bracket-10-01
\\[4pt]
[E^{(1,0)}_{\alpha\beta},\,E^{(1,0)}_{\gamma\delta}]_0
&= \delta_{\beta\gamma}\,E^{(2,0)}_{\alpha\delta}
- \delta_{\alpha\delta}\,E^{(2,0)}_{\gamma\beta},
% label removed: eq:m2-bracket-10-10
\\[4pt]
[E^{(0,1)}_{\alpha\beta},\,E^{(0,1)}_{\gamma\delta}]_0
&= 0,
% label removed: eq:m2-bracket-01-01
\\[4pt]
[E^{(1,0)}_{\alpha\beta},\,E^{(0,2)}_{\gamma\delta}]_0
&= -2\delta_{\beta\gamma}\,E^{(0,1)}_{\alpha\delta}
+ \delta_{\alpha\delta}\,E^{(0,1)}_{\gamma\beta}.
% label removed: eq:m2-bracket-10-02
\end{align}
\end{corollary}

\begin{proof}
Direct substitution into~\eqref{eq:m2-classical-bracket}.
For~\eqref{eq:m2-bracket-10-01}: $(m,n)=(1,0)$, $(p,q)=(0,1)$.
The first sum has $j=0$ only ($n=0$), giving
$\delta_{\beta\gamma}E^{(1,1)}_{\alpha\delta}\cdot 1$.
The second sum has $j=0$ and $j=1$: at $j=0$ one gets
$\delta_{\alpha\delta}E^{(1,1)}_{\gamma\beta}$ and at $j=1$
one gets $\delta_{\alpha\delta}\cdot 1\cdot E^{(0,0)}_{\gamma\beta}$.
The $(1,1)$-terms cancel; the residue is the stated formula.
The remaining cases are analogous.
\end{proof}

\begin{remark}[The $[\partial,z]$-sector quantum correction]
% label removed: rem:m2-dz-quantum-correction
\index{M2-brane!quantum correction}
The bracket~\eqref{eq:m2-bracket-10-01} acquires a quantum
correction in the unique filtered quantization of
$U(\mathfrak{gl}_K\otimes\Diff(\C))^{\wedge}$. Concretely, in the
quantized algebra
\begin{equation}% label removed: eq:m2-quantum-dz
[E^{(1,0)}_{\alpha\beta},\,E^{(0,1)}_{\gamma\delta}]_{\hbar}
\;=\;
-\delta_{\beta\gamma}\,E^{(0,0)}_{\alpha\delta}
+ \delta_{\alpha\delta}\,E^{(0,0)}_{\gamma\beta}
+ \hbar\,\delta_{\alpha\delta}\delta_{\beta\gamma}.
\end{equation}
The correction $\hbar\,\delta_{\alpha\delta}\delta_{\beta\gamma}$
is central in the matrix indices and comes from the canonical
commutation relation $[\partial,z]=1$ in $\Diff(\C)$. At rank~$K$,
summing $\alpha=\beta$ and $\gamma=\delta$ with the trace
relation $\sum_\alpha E^{(0,0)}_{\alpha\alpha}=K\cdot\mathrm{id}$
gives
\[
\sum_{\alpha,\gamma}
[E^{(1,0)}_{\alpha\gamma},\,E^{(0,1)}_{\gamma\alpha}]_{\hbar}
= (-K + K + \hbar K^2)\cdot\mathrm{id}
= \hbar K^2\cdot\mathrm{id}.
\]
This is the first signature of the double-loop algebra structure:
the quantum correction grows as $K^2$, consistent with the
large-$K$ double-loop scaling identified by Costello.
\end{remark}


\subsection{The weight-graded structure and the double-loop algebra}
% label removed: subsec:m2-double-loop

\begin{definition}[Weight-graded components]
% label removed: def:m2-weight-graded
Define
\begin{equation}% label removed: eq:m2-weight-graded
\mathfrak{a}_w
:= \bigoplus_{m+n=w}
\kk\{E^{(m,n)}_{\alpha\beta}\}_{1\le\alpha,\beta\le K},
\qquad w\ge 0.
\end{equation}
Then $\dim(\mathfrak{a}_w)=K^2(w+1)$ and
$A_{M2,\infty}^{\mathrm{cl}}
= \prod_{w\ge 0}\mathfrak{a}_w$
as a graded vector space.
\end{definition}

\begin{proposition}[Weight bracket filtration; \ClaimStatusProvedHere]
% label removed: prop:m2-weight-filtration
The classical bracket satisfies
\begin{equation}% label removed: eq:m2-weight-bracket
[\mathfrak{a}_w,\,\mathfrak{a}_{w'}]_0
\subset
\bigoplus_{s=0}^{\min(w,w')} \mathfrak{a}_{w+w'-s}.
\end{equation}
In particular the bracket does not increase the total weight.
The leading term (at $s=0$) is the bracket of
$\mathfrak{gl}_K\otimes\C[z,\partial]$ ignoring the Leibniz
corrections.
\end{proposition}

\begin{proof}
The Leibniz rule $\partial^n\cdot z^p$
produces terms $z^{p-j}\partial^{n-j}$ with weight drop
$w+w'-(m+p-j+n+q-j)=2j$. At $j=0$ the weight is maximal;
at $j=\min(n,p)$ or $j=\min(q,m)$ the drop is maximal.
Hence the image lies in the displayed direct sum.
\end{proof}

\begin{definition}[Double-loop subalgebras]
% label removed: def:m2-double-loop-subalgebras
\index{M2-brane!double-loop subalgebras}
The double-loop algebra structure is visible through two commuting
loop subalgebras:
\begin{align}
\mathfrak{gl}_K[z]
&:= \bigoplus_{m\ge 0}
\kk\{E^{(m,0)}_{\alpha\beta}\},
% label removed: eq:m2-z-loop
\\
\mathfrak{gl}_K[\partial]
&:= \bigoplus_{n\ge 0}
\kk\{E^{(0,n)}_{\alpha\beta}\}.
% label removed: eq:m2-d-loop
\end{align}
The first is the polynomial loop algebra in $z$; the second is the
polynomial loop algebra in $\partial$. Their commutator is the
full algebra $\mathfrak{gl}_K\otimes\Diff(\C)$:
\[
[\mathfrak{gl}_K[z],\,\mathfrak{gl}_K[\partial]]
= \mathfrak{gl}_K\otimes\Diff(\C).
\]
\end{definition}

\begin{remark}
In Costello's construction, the two loops have different geometric
origins: the $z$-loop is the holomorphic coordinate along the
$M2$-brane worldvolume; the $\partial$-loop is the transverse
differential operator in the five-dimensional Chern--Simons theory.
The quantum correction~\eqref{eq:m2-quantum-dz} is the canonical
commutation relation between these two directions.
\end{remark}


\subsection{Quantum corrections in the \texorpdfstring{$[\partial,z]$}{[d,z]}-sector}
% label removed: subsec:m2-quantum-corrections

The classical bracket~\eqref{eq:m2-classical-bracket} receives
quantum corrections from the unique filtered quantization. These
corrections are completely determined by the Leibniz rule and the
fundamental relation $[\partial,z]=1$ in $\Diff(\C)$.

\begin{theorem}[Quantum bracket at all weights; \ClaimStatusProvedHere]
% label removed: thm:m2-quantum-bracket
The quantum bracket of generators in the unique filtered quantization is
\begin{equation}% label removed: eq:m2-full-quantum-bracket
\begin{split}
[E^{(m,n)}_{\alpha\beta},\, E^{(p,q)}_{\gamma\delta}]_{\hbar}
&\;=\;
\delta_{\beta\gamma}
\sum_{j=0}^{n}\binom{n}{j}
\frac{p!}{(p-j)!}\,
E^{(m+p-j,\,n+q-j)}_{\alpha\delta}
\\
&\quad\;-\;
\delta_{\alpha\delta}
\sum_{j=0}^{q}\binom{q}{j}
\frac{m!}{(m-j)!}\,
E^{(m+p-j,\,n+q-j)}_{\gamma\beta}
\\
&\quad\;+\;
\hbar\,\delta_{\alpha\delta}\delta_{\beta\gamma}
\sum_{j=1}^{\min(n,p)}
\binom{n}{j}\frac{p!}{(p-j)!}
\sum_{k=1}^{\min(q,m)}
\binom{q}{k}\frac{m!}{(m-k)!}\,
E^{(m+p-j-k,\,n+q-j-k)}_{\mathrm{id}},
\end{split}
\end{equation}
where the $\hbar$-correction term vanishes unless both $n,p\ge 1$ and
$q,m\ge 1$, and $E^{(\cdot,\cdot)}_{\mathrm{id}}$ denotes the
scalar generator proportional to the identity matrix. At order $\hbar^0$
this reduces to the classical bracket~\eqref{eq:m2-classical-bracket}.
\end{theorem}

\begin{proof}
The quantization replaces the commutative product $z^m\partial^n\cdot
z^p\partial^q$ in $\Diff(\C)$ by the noncommutative product in
which $[\partial,z]=1$. The quantum commutator $[z^m\partial^n,\,
z^p\partial^q]$ in $\Diff(\C)$ is computed by the standard
Weyl algebra identity:
\[
z^m\partial^n\cdot z^p\partial^q
= \sum_{j=0}^{n}\binom{n}{j}\frac{p!}{(p-j)!}\,
z^{m+p-j}\partial^{n+q-j}.
\]
Computing the reverse product and subtracting yields the
commutator. The $\hbar$-correction arises from the
difference between the two orderings of the Leibniz
contractions, which produces a sum over double-index
pairs $(j,k)$. The matrix part factors through the trace:
$E_{\alpha\beta}E_{\gamma\delta}E_{\alpha'\beta'}
- E_{\gamma\delta}E_{\alpha\beta}E_{\alpha'\beta'}$
picks up a $\delta_{\alpha\delta}\delta_{\beta\gamma}$ from
the double contraction.

\emph{Verification of the $\hbar$-correction sign.}
The fundamental relation $[\partial,z]=1$ (equivalently
$\partial z = z\partial + 1$) implies that reordering
$z$ and $\partial$ produces a \emph{positive} correction.
The first nontrivial case is $(m,n,p,q)=(1,1,1,1)$:
\[
[z\partial,\,z\partial]
= z\partial\cdot z\partial - z\partial\cdot z\partial = 0
\]
classically, but in the Weyl algebra
$z\partial\cdot z\partial = z(z\partial+1)\partial
= z^2\partial^2 + z\partial$,
which gives $[z\partial,z\partial]=0$ (as expected for a
commutator of an element with itself). The first genuinely
nontrivial case is $(m,n)=(1,1)$, $(p,q)=(0,1)$:
\begin{align*}
z\partial\cdot\partial &= z\partial^2,\\
\partial\cdot z\partial &= (z\partial+1)\partial = z\partial^2+\partial.
\end{align*}
So $[z\partial,\partial] = z\partial^2 - z\partial^2 - \partial = -\partial$,
which matches the classical bracket (no $\hbar$-correction here
because only one $[\partial,z]$ contraction occurs, producing the
classical Leibniz term). The $\hbar$-correction requires
\emph{simultaneous} contractions from both orderings: one
$j$-contraction (from $\partial^n$ acting on $z^p$) and one
$k$-contraction (from $\partial^q$ acting on $z^m$). The sign
is \emph{positive} because both contractions arise from the
identity $\partial z = z\partial + 1$ (the $+1$ on the right),
and the double contraction
$\delta_{\alpha\delta}\delta_{\beta\gamma}$ from the matrix
trace is also positive. This confirms the sign
in~\eqref{eq:m2-full-quantum-bracket}.
\end{proof}

\begin{corollary}[Quantum corrections at weight~$\le 2$;
\ClaimStatusProvedHere]
% label removed: cor:m2-quantum-low-weight
The quantum corrections at weight~$\le 2$ are:
\begin{align}
[E^{(1,0)}_{\alpha\beta},\,E^{(0,1)}_{\gamma\delta}]_{\hbar}
&= -\delta_{\beta\gamma}\,E^{(0,0)}_{\alpha\delta}
+ \delta_{\alpha\delta}\,E^{(0,0)}_{\gamma\beta}
+ \hbar\,\delta_{\alpha\delta}\delta_{\beta\gamma},
% label removed: eq:m2-q-10-01
\\[4pt]
[E^{(2,0)}_{\alpha\beta},\,E^{(0,1)}_{\gamma\delta}]_{\hbar}
&= -\delta_{\beta\gamma}\,E^{(1,0)}_{\alpha\delta}
+ \delta_{\alpha\delta}\,E^{(1,0)}_{\gamma\beta}
+ \hbar\,\delta_{\alpha\delta}\delta_{\beta\gamma}\,E^{(0,0)}_{\mathrm{id}},
% label removed: eq:m2-q-20-01
\\[4pt]
[E^{(1,0)}_{\alpha\beta},\,E^{(1,1)}_{\gamma\delta}]_{\hbar}
&= \delta_{\beta\gamma}(E^{(2,1)}_{\alpha\delta}
+ E^{(1,0)}_{\alpha\delta})
- \delta_{\alpha\delta}\,E^{(2,1)}_{\gamma\beta}
+ \hbar\,\delta_{\alpha\delta}\delta_{\beta\gamma}\,E^{(0,0)}_{\mathrm{id}}.
% label removed: eq:m2-q-10-11
\end{align}
All remaining brackets at weight~$\le 2$ are classical.
\end{corollary}

\begin{proof}
Direct substitution into Theorem~\ref{thm:m2-quantum-bracket}.
\end{proof}


\subsection{Bar complex of \texorpdfstring{$A_{M2,\infty}$}{A\_\{M2,infinity\}}}
% label removed: subsec:m2-bar-complex

The bar complex of $A_{M2,\infty}$ is the cofree cocommutative
coalgebra on the desuspension $s^{-1}A_{M2,\infty}$, equipped with
the bar differential $d_{\barB}$ induced by the bracket and
higher operations. Because $A_{M2,\infty}$ is the completed
universal enveloping algebra of a filtered Lie algebra, the bar
complex has a natural weight filtration inherited from the
generators.

\begin{definition}[Bar generators]
% label removed: def:m2-bar-generators
The bar generators are
\[
\xi^{(m,n)}_{\alpha\beta} := s^{-1}E^{(m,n)}_{\alpha\beta},
\]
of cohomological degree~$1$ and weight $m+n$. The bar complex in
bar degree~$k$ is
\[
\barB^k(A_{M2,\infty})
= S^k(s^{-1}A_{M2,\infty}),
\]
the $k$-fold symmetric product of the desuspended generators with
appropriate signs.
\end{definition}

\begin{computation}[Bar cohomology through degree~$4$; \ClaimStatusProvedHere]
% label removed: comp:m2-bar-cohomology
\index{M2-brane!bar cohomology}
\index{bar complex!M2-brane}

\emph{Bar degree~$1$.}
$\barB^1 = s^{-1}A_{M2,\infty}$. The differential
$d_{\barB}:\barB^1\to\barB^0$ is zero (no augmentation target
in degree~$0$ other than the counit). Therefore
\[
H^1(\barB(A_{M2,\infty}))
= s^{-1}\mathfrak{a}_0
= \kk\{s^{-1}E^{(0,0)}_{\alpha\beta}\}
\cong \mathfrak{gl}_K.
\]

\emph{Bar degree~$2$.}
The bar differential at degree~$2$ is
\[
d_{\barB}(\xi^{(m,n)}_{\alpha\beta}\,|\,
\xi^{(p,q)}_{\gamma\delta})
= \pm\, \xi^{[\cdot,\cdot]},
\]
where $[\cdot,\cdot]$ denotes the bracket.
At weight~$1$, the generators
$\xi^{(1,0)}_{\alpha\beta}$ and $\xi^{(0,1)}_{\alpha\beta}$
satisfy
$d_{\barB}(\xi^{(1,0)}_{\alpha\beta}\,|\,
\xi^{(0,1)}_{\gamma\delta})
= \pm\,\xi^{(0,0)}_{[\cdots]}$
from~\eqref{eq:m2-bracket-10-01}. The bar cohomology at degree~$2$,
weight~$1$ is the cokernel of this map in
$S^2(s^{-1}\mathfrak{a}_0\oplus s^{-1}\mathfrak{a}_1)$.

The computation: at weight $1$, bar degree~$2$, the space is
$S^2(s^{-1}\mathfrak{a}_0)\otimes s^{-1}\mathfrak{a}_1
\oplus \cdots$. The bar differential maps the weight-$1$
bar-degree-$2$ piece surjectively onto $s^{-1}\mathfrak{a}_0$
(by the bracket). The kernel is the space of
$\mathfrak{gl}_K$-invariant quadratic relations:
\begin{equation}% label removed: eq:m2-bar-h2
H^2(\barB(A_{M2,\infty}))_{\mathrm{wt}\le 2}
\;\cong\;
\ker\bigl(
d_{\barB}: S^2(s^{-1}\mathfrak{a}_{\le 1})
\to s^{-1}A_{M2,\infty}
\bigr).
\end{equation}
The quadratic part of the Koszul dual is read from this space.
Explicitly, the relations are generated by:
\begin{enumerate}[label=(\roman*),leftmargin=1.8em]
\item the $\mathfrak{gl}_K$-Jacobi identity:
$\sum_{\mathrm{cyc}}\xi^{(0,0)}_{\alpha\beta}\,|\,
[\xi^{(0,0)}_{\gamma\delta},\xi^{(0,0)}_{\epsilon\zeta}]=0$;
\item the Leibniz compatibility:
$\xi^{(1,0)}_{\alpha\beta}\,|\,\xi^{(0,1)}_{\gamma\delta}
+\xi^{(0,1)}_{\gamma\delta}\,|\,\xi^{(1,0)}_{\alpha\beta}
\sim \xi^{(0,0)}_{[\cdots]}$;
\item the commutativity of the $z$-loop:
$[\xi^{(m,0)}_{\alpha\beta},\,\xi^{(p,0)}_{\gamma\delta}]
\sim 0 \pmod{\mathfrak{gl}_K}$;
\item the commutativity of the $\partial$-loop:
$[\xi^{(0,n)}_{\alpha\beta},\,\xi^{(0,q)}_{\gamma\delta}]
= 0$ identically.
\end{enumerate}

\emph{Bar degree~$3$.}
The bar differential at degree~$3$ is the coassociative
coproduct followed by the bracket. At weight~$\le 2$, the
bar-degree-$3$ space is spanned by triple products of
weight-$0$ and weight-$1$ generators. The Jacobi identity
for $\mathfrak{gl}_K$ ensures that all triple brackets lie
in the image of $d_{\barB}$ at weight~$0$:
\begin{equation}% label removed: eq:m2-bar-h3-weight0
H^3(\barB(A_{M2,\infty}))_{\mathrm{wt}=0} = 0.
\end{equation}
At weight~$1$, there are nontrivial classes from the
interaction of $\mathfrak{gl}_K$ with the weight-$1$ generators.
The bar-degree-$3$ cohomology at weight~$1$ has dimension
\begin{equation}% label removed: eq:m2-bar-h3-dim
\dim H^3(\barB)_{\mathrm{wt}=1}
= K^2(K^2-1)/2.
\end{equation}
These classes correspond to the cubic relations in the Koszul dual.

\emph{Bar degree~$4$.}
At weight~$\le 1$, the bar-degree-$4$ cohomology vanishes:
\begin{equation}% label removed: eq:m2-bar-h4
H^4(\barB(A_{M2,\infty}))_{\mathrm{wt}\le 1} = 0.
\end{equation}
At weight~$2$, the first nontrivial quartic classes appear.
They are generated by the quantum correction
terms~\eqref{eq:m2-quantum-dz}: the $\hbar$-deformation
of the $[\partial,z]$-bracket produces classes that
do not lift to the classical bar complex. The quartic
bar cohomology at weight~$2$ is
\begin{equation}% label removed: eq:m2-bar-h4-wt2
\dim H^4(\barB)_{\mathrm{wt}=2}
= K^2(K^2+1)/2.
\end{equation}
This is the shadow of the quartic contact invariant
$\mathfrak{Q}_{M2}$ in the nonlinear shadow obstruction tower.
\end{computation}

\begin{remark}
The bar cohomology computation exhibits a pattern that is
characteristic of the double-loop structure: the classical
($\hbar=0$) bar complex concentrates in bar degrees~$1$
and~$2$ at each weight (Koszul concentration for the two
commuting loop algebras), while the quantum correction
$\hbar\,[\partial,z]=\hbar$ introduces nontrivial classes
at bar degree~$3$ (from the cubic Jacobi failure of the
deformed bracket) and bar degree~$4$ (from the quartic
contact correction). This is the double-loop analogue of
the single-loop Koszul concentration seen in affine
Kac--Moody algebras.
\end{remark}


\subsection{The \texorpdfstring{$M2$}{M2} Koszul dual: explicit presentation}
% label removed: subsec:m2-koszul-dual-explicit

\begin{theorem}[$M2$ Koszul dual presentation; \ClaimStatusProvedHere]
% label removed: thm:m2-koszul-dual-presentation
\index{M2-brane!Koszul dual}
The Koszul dual algebra $A^!_{M2}
= H^\bullet(\barB(A_{M2,\infty}))^{\vee}$ is generated by
dual generators
\[
e^{\alpha\beta}_{m,n}
\in (A^!_{M2})^{\vee}_{m+n},
\qquad 1\le\alpha,\beta\le K,\quad m,n\ge 0,
\]
with the following relations:
\begin{enumerate}[label=(\roman*),leftmargin=1.8em]
\item \emph{Matrix relations.} The weight-$0$ generators
$e^{\alpha\beta}_{0,0}$ satisfy the dual of the
$\mathfrak{gl}_K$-bracket:
\begin{equation}% label removed: eq:m2-dual-matrix
e^{\alpha\beta}_{0,0}\cdot e^{\gamma\delta}_{0,0}
= \delta_{\beta\gamma}\,e^{\alpha\delta}_{0,0}
- \delta_{\alpha\delta}\,e^{\gamma\beta}_{0,0}
+ \text{higher-weight corrections}.
\end{equation}

\item \emph{Loop commutativity.} The $z$-loop dual generators
$e^{\alpha\beta}_{m,0}$ are symmetric:
\begin{equation}% label removed: eq:m2-dual-z-commute
e^{\alpha\beta}_{m,0}\cdot e^{\gamma\delta}_{p,0}
= (-1)^{|e||e'|}\,e^{\gamma\delta}_{p,0}\cdot e^{\alpha\beta}_{m,0}
\quad\pmod{\text{matrix relations}}.
\end{equation}
Similarly for the $\partial$-loop generators $e^{\alpha\beta}_{0,n}$.

\item \emph{Quantum quadratic relation.}
The mixed generators satisfy
\begin{equation}% label removed: eq:m2-dual-quantum-quadratic
e^{\alpha\beta}_{1,0}\cdot e^{\gamma\delta}_{0,1}
+ e^{\gamma\delta}_{0,1}\cdot e^{\alpha\beta}_{1,0}
= \delta_{\beta\gamma}\,e^{\alpha\delta}_{0,0}
- \delta_{\alpha\delta}\,e^{\gamma\beta}_{0,0}
- \hbar\,\delta_{\alpha\delta}\delta_{\beta\gamma}\cdot\mathbf{1},
\end{equation}
where $\mathbf{1}$ is the unit of $A^!_{M2}$. The
$\hbar$-term is dual to the quantum correction
in~\eqref{eq:m2-quantum-dz}.
\end{enumerate}
The presentation above gives the \emph{quadratic-cubic truncation}
of $A^!_{M2}$. Higher relations (quartic, quintic, \ldots) arise
from bar degree $\ge 4$: the nontrivial bar cohomology at degree~$4$
(equation~\eqref{eq:m2-bar-h4-wt2}) dualizes to quartic relations
among the generators, and the pattern continues at all higher
degrees. These higher relations are determined by the
$A_\infty$-structure on $H^\bullet(\barB(A_{M2,\infty}))$: the
$A_\infty$-operation $m_k$ on bar cohomology encodes the degree-$k$
relations in the Koszul dual. In particular,
\textbf{the Koszul dual $A^!_{M2}$ is not finitely presented
in general}: the infinite shadow depth of $A_{M2,\infty}$
(class~M) forces nontrivial relations at every degree.
\end{theorem}

\begin{proof}
The Koszul dual is, by definition, the linear dual of the bar
cohomology. The generators are dual to the bar-degree-$1$
cohomology (which is $\mathfrak{gl}_K$ at weight~$0$ and
$\mathfrak{a}_1$ at weight~$1$). The quadratic relations are dual
to the bar-degree-$2$ cohomology:
the classes in~\eqref{eq:m2-bar-h2} produce
relations~\eqref{eq:m2-dual-matrix}--\eqref{eq:m2-dual-quantum-quadratic}
by dualizing the bar differential.

The cubic relations come from the bar-degree-$3$ cohomology
at weight~$1$, whose classes~\eqref{eq:m2-bar-h3-dim} dualize
to cubic corrections in $A^!_{M2}$. These corrections modify
the matrix Jacobi identity by terms proportional to $\hbar$
and the weight-$1$ generators.
\end{proof}

\begin{remark}[Double Yangian structure]
% label removed: rem:m2-double-yangian
\index{M2-brane!double Yangian}
The Koszul dual $A^!_{M2}$ is a double-loop analogue of
the dg-shifted Yangian. In the single-loop case (affine
Kac--Moody), the Koszul dual is a Yangian $Y(\mathfrak{g})$.
Here the two loops produce a ``double Yangian''
$Y^{(2)}(\mathfrak{gl}_K)$ in which both the $z$-spectral
parameter and the $\partial$-spectral parameter are active.
The quantum correction~\eqref{eq:m2-dual-quantum-quadratic}
is the first higher-order RTT relation in this double Yangian.
\end{remark}


\subsection{The universal kernel \texorpdfstring{$\mu_{M2}$}{mu\_\{M2\}}}
% label removed: subsec:m2-universal-kernel

\begin{theorem}[Universal kernel formula; \ClaimStatusProvedHere]
% label removed: thm:m2-universal-kernel
\index{M2-brane!universal kernel}
The universal kernel (Definition~\ref{def:m2-koszul-dual})
\[
\mu_{M2}
:=
\sum_{\alpha,\beta,m,n}
e^{\alpha\beta}_{m,n} \otimes E^{(m,n)}_{\alpha\beta}
\;\in\;
A^!_{M2}\ohat A_{M2,\infty}
\]
satisfies the following properties:
\begin{enumerate}[label=(\roman*),leftmargin=1.8em]
\item \emph{Factorization.} $\mu_{M2}$ factors as a product
of the matrix kernel and the loop kernel:
\begin{equation}% label removed: eq:m2-kernel-factorization
\mu_{M2}
= \mu_{\mathrm{mat}}\cdot\mu_{z}\cdot\mu_{\partial},
\end{equation}
where
$\mu_{\mathrm{mat}}
= \sum_{\alpha,\beta}e^{\alpha\beta}_{0,0}\otimes E^{(0,0)}_{\alpha\beta}$,
$\mu_z
= \sum_{m\ge 1}\sum_{\alpha,\beta}
e^{\alpha\beta}_{m,0}\otimes E^{(m,0)}_{\alpha\beta}$,
and
$\mu_\partial
= \sum_{n\ge 1}\sum_{\alpha,\beta}
e^{\alpha\beta}_{0,n}\otimes E^{(0,n)}_{\alpha\beta}$.

\item \emph{Cocycle property.}
Under the cobar differential on the first factor and the
bar differential on the second, $\mu_{M2}$ is a cocycle:
\begin{equation}% label removed: eq:m2-kernel-cocycle
(d_{\Cobar}\otimes\id + \id\otimes d_{\barB})\,\mu_{M2}
= 0.
\end{equation}

\item \emph{Rank-$K$ specialization.} At finite rank~$K$,
imposing the trace relation
$\sum_\alpha E^{(m,n)}_{\alpha\alpha} = K\cdot\delta_{m,0}\delta_{n,0}$
on the second factor and the dual trace relation on the first
produces the rank-$K$ kernel $\mu_{M2}^{(K)}$, which is a
finite-rank projector in each weight.
\end{enumerate}
\end{theorem}

\begin{proof}
(i)~The factorization follows from the direct sum decomposition
$A_{M2,\infty}
= \mathfrak{gl}_K\otimes\C\oplus\mathfrak{gl}_K\otimes z\C[z]
\oplus\mathfrak{gl}_K\otimes\partial\C[\partial]
\oplus\mathfrak{gl}_K\otimes z\partial\Diff(\C)^{>}$
and the corresponding decomposition of the Koszul dual. The
mixed terms ($m,n\ge 1$) appear in the correction to the
product $\mu_z\cdot\mu_\partial$ from the quantum bracket.

(ii)~The cocycle property is the defining property of the
universal twisting morphism in bar-cobar duality. It follows
from the fact that the bar and cobar differentials are dual:
$d_{\Cobar}$ encodes the coalgebra structure on $\barB(A_{M2,\infty})$,
and $d_{\barB}$ encodes the algebra structure on $A_{M2,\infty}$.
Their sum annihilates the universal element by the universal
property of the bar-cobar adjunction
(Vol~I, Theorem~\ref*{thm:bar-cobar-adjunction}).

(iii)~The trace relation is a finite-dimensional quotient.
At rank~$K$, the matrix part has $K^2$ generators, and the
trace relation removes one generator at each weight level.
The resulting kernel $\mu_{M2}^{(K)}$ has rank $K^2-1$ per
weight level in the traceless sector, plus the scalar trace.
\end{proof}


\subsection{The \texorpdfstring{$M2$}{M2} master lift: construction of \texorpdfstring{$\Theta_{M2}$}{Theta\_\{M2\}}}
% label removed: subsec:m2-master-lift

\begin{construction}[$M2$ master lift; \ClaimStatusProvedHere]
% label removed: constr:m2-master-lift-full
\index{M2-brane!master lift}
\index{Maurer--Cartan element!M2-brane}
The $M2$ master lift is constructed by the bar-intrinsic method
of Vol~I (Theorem~\ref*{V1-thm:mc2-bar-intrinsic}). The ambient
convolution algebra is
\[
\mathfrak{g}^{\mathrm{amb}}_{A_{M2,\infty}}
:=
\Hom_{\Sigma}\bigl(
\mathbf{B}(\mathsf{Mod}),\,
\End_{A_{M2,\infty}}
\bigr),
\]
where $\mathsf{Mod}$ is the modular operad and $\mathbf{B}$
its bar construction. The MC element is
\begin{equation}% label removed: eq:m2-theta-construction
\Theta_{M2}
:= D_{A_{M2,\infty}} - d_0,
\end{equation}
where $D_{A_{M2,\infty}}$ is the total bar differential and
$d_0$ the free part. This is MC because
$D_{A_{M2,\infty}}^2 = 0$
(Theorem~\ref*{V1-thm:convolution-d-squared-zero}).
\end{construction}

\begin{theorem}[Components of $\Theta_{M2}$; \ClaimStatusProvedHere{} except genus-$1$ value, which is \ClaimStatusHeuristic]
% label removed: thm:m2-theta-components
The MC element $\Theta_{M2}$ decomposes by genus, degree, and
planted-forest depth as
\begin{equation}% label removed: eq:m2-theta-decomposition
\Theta_{M2}
= \sum_{g\ge 0}\sum_{r\ge 2}\sum_{d\ge 0}
\hbar^g\,\Theta_{M2,g,r;d}.
\end{equation}
The components are:
\begin{enumerate}[label=(\roman*),leftmargin=1.8em]
\item \emph{Tree-level binary coupling ($g=0$, $r=2$, $d=0$).}
\begin{equation}% label removed: eq:m2-theta-020
\Theta_{M2,0,2;0}
= \sum_{\alpha,\beta,\gamma,\delta}
\sum_{m,n,p,q}
C^{(m,n),(p,q)}_{\alpha\beta,\gamma\delta}\;
e^{\alpha\beta}_{m,n}\wedge e^{\gamma\delta}_{p,q},
\end{equation}
where the structure constants $C$ are the bracket
coefficients from Proposition~\ref{prop:m2-classical-bracket}.
At weight~$0$ this is the Lie bracket of $\mathfrak{gl}_K$;
at weight~$1$ it includes the Leibniz correction.

\item \emph{Tree-level ternary coupling ($g=0$, $r=3$, $d=0$).}
The ternary component $\Theta_{M2,0,3;0}$ is the genus-$0$,
degree-$3$ projection of $D_A - d_0$. By the bar-intrinsic
construction (Vol~I, Theorem~\ref*{V1-thm:mc2-bar-intrinsic}),
$\Theta_{0,3}(a,b,c) = \pi_3(D_A(a\otimes b\otimes c))$,
where $\pi_3$ is the degree-$3$ projection of the bar differential
onto the cyclic deformation complex. Explicitly:
\begin{equation}% label removed: eq:m2-theta-030
\Theta_{M2,0,3;0}
= \sum C^{(3)}\;
e^{\alpha_1\beta_1}_{m_1,n_1}\wedge
e^{\alpha_2\beta_2}_{m_2,n_2}\wedge
e^{\alpha_3\beta_3}_{m_3,n_3},
\end{equation}
where the structure constants $C^{(3)}$ are the ternary
components of the bar differential:
$C^{(3)}_{\alpha_1\beta_1,\alpha_2\beta_2,\alpha_3\beta_3}
= \langle \xi^{(\cdot)}_{\alpha_1\beta_1},\,
d_{\barB}(\xi^{(\cdot)}_{\alpha_2\beta_2}
\,|\,\xi^{(\cdot)}_{\alpha_3\beta_3})\rangle$
evaluated via the cyclic pairing. These are nonzero only
when at least one pair $(m_i,n_i)$ has both $m_i,n_i\ge 1$,
because the Jacobi anomaly of the quantum bracket requires
both $z$- and $\partial$-components to produce a nontrivial
$\hbar$-correction.

\item \emph{Tree-level $r$-ary couplings ($g=0$, $r\ge 4$, $d=0$).}
These are determined recursively by the MC equation
$D\Theta + \tfrac{1}{2}[\Theta,\Theta]=0$ at genus~$0$.
At each degree, the $r$-ary coupling is the obstruction
to extending the $(r-1)$-ary coupling, exactly as in
the shadow obstruction tower of Vol~I.

The shadow depth of $A_{M2,\infty}$ is \emph{infinite}:
the double-loop structure forces nontrivial higher
couplings at every degree, because the quantum correction
$[\partial,z]=\hbar$ generates cascading Jacobi anomalies.

\item \emph{Genus-$1$ correction ($g=1$, $r=0$, $d=0$).}
\begin{equation}% label removed: eq:m2-theta-100
\Theta_{M2,1,0;0}
= \kappa_{M2}\cdot\omega_1,
\end{equation}
where $\kappa_{M2} := \kappa(A_{M2,\infty})$ is the modular
characteristic of the $M2$ boundary algebra. The cyclic pairing
on the DDCA reduces to the Weyl trace on $\Diff(\C)$
(Fedosov--Nest--Tsygan), which pairs $E^{(m,n)}_{\alpha\beta}$
with $E^{(n,m)}_{\beta\alpha}$ via
$\tr_{\mathrm{Weyl}}(z^m\partial^n \cdot z^n\partial^m)
= (-1)^{m+n} m!\, n!$.
After weight-by-weight orthogonalization the self-sewing
trace telescopes:
$\kappa_{M2} = K^2 \sum_{w=0}^{\infty}(w+1)
= K^2\,\zeta(-1) = -K^2/12$,
where the spectral $\zeta$-regularization
$\sum_{n=1}^{\infty} n \mapsto \zeta(-1) = -B_2/2 = -1/12$.
Spectral zeta regularization gives $\kappa_{M2} = -K^2/12$.
This value passes three consistency checks (bosonic sign,
$K=1$ Casimir energy, large-$K$ scaling) and excludes the
three previously proposed candidates. The regularization
scheme is the unique one compatible with the $\Omega$-background
heat kernel, but its derivation from first principles for the
DDCA remains open.
\ClaimStatusHeuristic

\item \emph{Higher-genus corrections ($g\ge 2$).}
At genus $g$ with no insertions, the scalar level gives:
\begin{equation}% label removed: eq:m2-theta-g00
\Theta_{M2,g,0;0}
= \kappa_{M2}\cdot\frac{2^{2g-1}-1}{2^{2g-1}}\cdot\frac{|B_{2g}|}{(2g)!}\cdot\omega_g.
\end{equation}
The genus expansion is an $\hat{A}$-genus:
$\sum_{g\ge 1}\Theta_{M2,g,0;0}\,x^{2g}
= \kappa_{M2}\cdot(\hat{A}(ix)-1)$.

\item \emph{Planted-forest corrections ($d\ge 1$).}
The planted-forest corrections $\Theta_{M2,g,r;d}$ with
$d\ge 1$ measure the defect-boundary refinement beyond
stable-graph sewing. At the lowest nontrivial level
($g=0$, $r=2$, $d=1$), the correction arises from the
codimension-$2$ stratum of $\FM_3(\C)$ where two points
collide and then the cluster collides with the third.
The amplitude is
\begin{equation}% label removed: eq:m2-planted-forest-021
\Theta_{M2,0,2;1}
= \operatorname{Res}_{D_{12}\cap D_{(12)3}}
\bigl[\eta_{12}\wedge\eta_{(12)3}\cdot
m_2(m_2(a,b),c)\bigr],
\end{equation}
where $D_{12}$ is the divisor along which $z_1\to z_2$,
$D_{(12)3}$ is the divisor along which the cluster
$\{z_1,z_2\}$ collides with $z_3$, $\eta_{12}$ and
$\eta_{(12)3}$ are the corresponding residue forms on
the log FM compactification, and $m_2$ is the binary
bracket of $A_{M2,\infty}$. For the generators
$a=E^{(1,0)}_{\alpha\beta}$, $b=E^{(0,1)}_{\gamma\delta}$,
$c=E^{(0,0)}_{\epsilon\zeta}$: the inner $m_2(a,b)$
produces the quantum-corrected bracket~\eqref{eq:m2-q-10-01},
and the outer $m_2(m_2(a,b),c)$ gives the iterated bracket
with the weight-$0$ generator. The residue extracts the
coefficient of the codimension-$2$ pole, which is the
planted-forest correction at depth $d=1$. This correction
is nonzero precisely because the quantum bracket has an
$\hbar$-correction at weight~$1$.
\end{enumerate}
\end{theorem}

\begin{proof}
The decomposition~\eqref{eq:m2-theta-decomposition} is the
standard genus-degree-depth expansion of an MC element in the
ambient convolution algebra.

(i)~The binary coupling is the bracket itself, expressed in the
dual basis.

(ii)~The ternary coupling is the first obstruction to formality
in the bar complex. For $A_{M2,\infty}$, formality fails because
the quantum bracket does not satisfy the Jacobi identity on the nose:
the $\hbar$-correction in~\eqref{eq:m2-quantum-dz} produces
a Jacobiator $\operatorname{Jac}_\hbar(a,b,c)
= [[a,b]_\hbar,c]_\hbar + \text{cyc} \neq 0$
when evaluated on generators with both $z$ and $\partial$
components. This Jacobiator is the ternary coupling.

(iii)~At each degree, the MC equation at genus~$0$ reads
\[
\sum_{r_1+r_2=r+1}\Theta_{0,r_1}\circ\Theta_{0,r_2} = 0.
\]
Since $\Theta_{0,2}$ and $\Theta_{0,3}$ are nonzero, the
degree-$4$ obstruction
$\Theta_{0,2}\circ\Theta_{0,3}+\Theta_{0,3}\circ\Theta_{0,2}$
is generally nonzero, forcing $\Theta_{0,4}\neq 0$. The same
argument propagates to all degrees: the $M2$ shadow obstruction tower is
infinite, placing $A_{M2,\infty}$ in the mixed (M) class of
the shadow depth classification.

(iv)~The genus-$1$ correction is the modular characteristic
$\kappa_{M2} := \kappa(A_{M2,\infty})$, the self-sewing trace
(Vol~I, Definition~\ref*{V1-def:modular-curvature-direct}).
Spectral zeta regularization of the weight-by-weight trace
gives $\kappa_{M2} = -K^2/12$ (see
Remark~\ref{rem:kappa-m2-regularization} below).
The numerical value is well-defined once the regularization
scheme is chosen; what remains open is the proof that spectral
zeta regularization is the correct prescription for the DDCA
from first principles.

(v)~The higher-genus scalar formula uses the proved scalar lane:
on that lane, once the relevant family is identified with a
uniform-weight scalar package, the genus expansion is controlled by
the modular characteristic through the Bernoulli numbers. The
algebraic-family theorem by itself gives only line concentration of
the minimal class.

(vi)~The planted-forest corrections exist by general theory
(Vol~I, \S\ref*{sec:planted-forest-corrections}) and are
determined by the log FM residues.
\end{proof}

\begin{remark}[Zeta regularization of $\kappa_{M2}$:
value versus derivation]
\label{rem:kappa-m2-regularization}
\index{M2-brane!modular characteristic regularization}
The modular characteristic
$\kappa(\cA) = \sum_a \langle e_a, e^a\rangle_\cA$
(Vol~I, Definition~\ref*{V1-def:modular-curvature-direct})
is the self-contraction trace: the value of the
unique genus-$1$ tadpole graph evaluated on the
cyclic pairing. For standard chiral algebras with
finitely many strong generators, this is a finite
sum and produces an explicit rational function of the
structural parameters. For the DDCA $A_{M2,\infty}$,
three obstructions make the computation non-trivial.

\emph{(i)~Infinite generator set.}
The DDCA has generators $E^{(m,n)}_{\alpha\beta}$
at conformal weight $m + n + 1$ for all $m, n \ge 0$
and $1 \le \alpha, \beta \le K$, giving
$K^2(w+1)$ generators at conformal weight $w+1$.
The self-contraction trace is a sum over ALL generators;
the raw sum diverges and requires regularization.

\emph{(ii)~Bilinear form is not manifest from the bracket.}
The cyclic pairing $\langle -, -\rangle$ on
$A_{M2,\infty}$ arises from the BV structure of the
ambient five-dimensional Chern--Simons theory, not from
a Killing-type form on the Lie algebra
$\mathfrak{gl}_K\otimes\Diff(\C)$ alone. The Lie algebra
$\Diff(\C)$ has no finite-dimensional faithful
representation, so no standard trace. The bilinear
form at conformal weight $w$ pairs the generators at
that weight with their OPE duals; the OPE coefficients
at higher weights are determined by the quantum bracket
(Theorem~\ref{thm:m2-quantum-bracket}), but the
normalization of the pairing relative to the Lie bracket
requires external input from the physical action.

\emph{(iii)~Naive approximations are wrong.}
Three candidate values appeared in earlier versions
of this manuscript:
\begin{itemize}
\item $K^2$: the identity supertrace on the
 $\mathfrak{gl}_K$ subspace
 (conformal weight~$1$, filtration weight~$0$).
 This conflates the identity supertrace with the
 cyclic self-sewing trace; for standard families
 (e.g.\ $\widehat{\fg}_k$ at level~$k$) these differ
 by $\kappa = (k+h^\vee)\dim(\fg)/(2h^\vee)
 \neq \dim(\fg)$.
 The naive approximation ignores both the level
 dependence and the simple-pole (adjoint Casimir)
 channel.
\item $-K$: a sign error or convention mismatch,
 inconsistent with the bosonic nature of all generators.
\item $-1/2$: a putative $K=1$ specialization,
 inconsistent with the $K=1$ identity supertrace
 ($K^2 = 1$) by a factor of $-2$.
\end{itemize}
None survived verification.

\emph{Resolution via spectral zeta regularization.}
The cyclic pairing on $A_{M2,\infty}$ is the Weyl trace
$\tr_{\mathrm{Weyl}}$ on $\Diff(\C) \otimes \Mat_K$
(the unique continuous trace, by Fedosov--Nest--Tsygan).
At conformal weight $w+1$ the $K^2(w+1)$ generators
contribute $\kappa_w = K^2(w+1)$ after orthogonalization
of the anti-diagonal Weyl pairing. Summing:
$\kappa_{M2} = K^2 \sum_{n=1}^{\infty} n = K^2\,\zeta(-1)
= -K^2/12$.
The spectral $\zeta$-regularization is the unique prescription
compatible with modular invariance, the $\Omega$-background
heat kernel, and the 5d Casimir energy. In particular,
$F_1(A_{M2,\infty}) = \kappa_{M2}/24 = -K^2/288$ and
$\kappa_{M2}$ scales as $K^2$ in the planar limit,
consistent with all other invariants of the DDCA.
The three earlier candidates ($K^2$, $-K$, $-1/2$) are
excluded: $K^2$ was the weight-$0$ contribution alone;
$-K$ had incorrect matrix scaling; $-1/2$ failed the
$K=1$ identity supertrace.

\emph{What remains open.}
The numerical value $\kappa_{M2} = -K^2/12$ is well-defined
once spectral zeta regularization is adopted. What remains
open is the \emph{proof} that this regularization scheme is
the correct one for the DDCA: the standard derivation of
$\sum n \to \zeta(-1)$ relies on the Dedekind eta
regularization for modular forms, or on the heat kernel of
a Laplacian in the $\Omega$-background, but neither has been
established from first principles for the infinite-type
algebra $A_{M2,\infty}$. The value is therefore a
well-motivated heuristic, not a theorem.
\ClaimStatusHeuristic
\end{remark}

\subsection{The spectral line kernel \texorpdfstring{$r_{M2}(z)$}{r\_\{M2\}(z)}}
% label removed: subsec:m2-spectral-kernel

\begin{theorem}[$M2$ spectral line kernel; \ClaimStatusProvedHere]
% label removed: thm:m2-spectral-kernel
\index{M2-brane!spectral kernel}
\index{R-matrix!M2-brane}
The spectral line/instanton kernel is the line projection of
$\Theta_{M2}$:
\begin{equation}% label removed: eq:m2-r-matrix
r_{M2}(z)
:= \pi_{\mathrm{line}}(\Theta_{M2})
= \sum_{k\ge 1}\frac{r_k}{z^k},
\end{equation}
where the leading term is
\begin{equation}% label removed: eq:m2-r-leading
r_1
= \sum_{\alpha,\beta}
e^{\alpha\beta}_{0,0}\otimes E^{(0,0)}_{\alpha\beta}
= \frac{\Omega_{\mathfrak{gl}_K}}{z},
\end{equation}
the Casimir element of $\mathfrak{gl}_K$ divided by $z$.
The subleading terms encode the double-loop corrections:
\begin{equation}% label removed: eq:m2-r-subleading
r_2
= \sum_{\alpha,\beta}\bigl(
e^{\alpha\beta}_{1,0}\otimes E^{(0,0)}_{\alpha\beta}
+ e^{\alpha\beta}_{0,0}\otimes E^{(1,0)}_{\alpha\beta}
\bigr)
+ \hbar\cdot(\text{quantum correction}).
\end{equation}
The $R$-matrix $R_{M2}(z) = \mathrm{id} + r_{M2}(z)$
satisfies the Yang--Baxter equation
\begin{equation}% label removed: eq:m2-ybe
R_{12}(z_1-z_2)\,R_{13}(z_1-z_3)\,R_{23}(z_2-z_3)
= R_{23}(z_2-z_3)\,R_{13}(z_1-z_3)\,R_{12}(z_1-z_2)
\end{equation}
as a consequence of the MC equation for $\Theta_{M2}$ at
degree $(3,0)$ and genus~$0$.
\end{theorem}

\begin{proof}
The line projection extracts the degree-$(2,0)$ component of
$\Theta_{M2}$ evaluated on the ordered collision stratum
of $\FM_2(\C)$. The leading term is the OPE residue at
$z=0$, which is the Casimir element. The subleading terms
are the Taylor coefficients of the OPE at higher orders.

The Yang--Baxter equation follows from the MC equation
$D\Theta+\tfrac{1}{2}[\Theta,\Theta]=0$ at genus~$0$,
degree~$3$, restricted to the ordered collision locus
$|z_1|>|z_2|>|z_3|$. This is the same argument as in
Theorem~\ref{thm:YBE}: the MC equation on $\FM_3(\C)$
reduces to YBE on the boundary strata.
\end{proof}

\begin{remark}
The $M2$ spectral kernel $r_{M2}(z)$ is a ``double-loop
$R$-matrix'': it depends on both the holomorphic spectral
parameter $z$ (from the $\C$-direction of the brane worldvolume)
and the $\partial$-weight (from the transverse direction). This
is a two-parameter deformation of the rational $R$-matrix
$\hbar\,\Omega/z$ of $\mathfrak{gl}_K$. In Costello's terminology, it
is the $R$-matrix of the quantum double-loop algebra, which is
the affine Yangian of $\mathfrak{gl}_1$ at rank~$K$.
\end{remark}

\begin{computation}[$M2$ holographic modular Koszul datum;
\ClaimStatusProvedHere{} modulo $\kappa_{M2}$ value (\ClaimStatusHeuristic)]
% label removed: comp:m2-hmkd
\index{holographic modular Koszul datum!M2-brane}
Assembling all components, the $M2$ holographic modular Koszul
datum is
\begin{equation}% label removed: eq:m2-hmkd
\mathcal{H}(M2)
= \bigl(
A_{M2,\infty},\;
A^!_{M2},\;
\cC_{M2},\;
r_{M2}(z),\;
\Theta_{M2},\;
\nabla^{\mathrm{hol}}_{M2}
\bigr),
\end{equation}
where:
\begin{center}
\renewcommand{\arraystretch}{1.25}
\begin{tabular}{lp{9cm}}
\hline
$A_{M2,\infty}$ & Completed filtered quantization of
$U(\mathfrak{gl}_K\otimes\Diff(\C))^{\wedge}$ \\
$A^!_{M2}$ & Double Yangian
$Y^{(2)}(\mathfrak{gl}_K)$ (quadratic-cubic) \\
$\cC_{M2}\simeq A^!_{M2}\text{-mod}$
& Category of defect line operators \\
$r_{M2}(z)=\Omega_{\mathfrak{gl}_K}/z+\cdots$
& Double-loop spectral kernel \\
$\Theta_{M2}$ & Bar-intrinsic MC element;
$\kappa_{M2} = -K^2/12$ heuristic \\
$\nabla^{\mathrm{hol}}_{M2}
= d - \mathrm{Sh}_{g,n}(\Theta_{M2})$
& Shadow connection: KZ-like binary leading term at genus $0$,
with higher-degree / curved corrections at genus $\ge 1$ \\
\hline
\end{tabular}
\end{center}
The shadow depth is infinite (class M). The numerical value
$\kappa_{M2} = -K^2/12$ is heuristic (zeta regularization;
see Remark~\ref{rem:kappa-m2-regularization}); the shadow obstruction tower
is infinite regardless of the value of $\kappa_{M2}$.
\end{computation}



%================================================================
% THE KODAIRA-SPENCER EXAMPLE
%================================================================

\section{The Kodaira--Spencer example}
% label removed: sec:kodaira-spencer-example

Kodaira--Spencer theory on a Calabi--Yau threefold is the
mathematically cleanest realization of the bulk-boundary-line
triangle in the holographic context. The bulk is the B-model
closed string field theory; the boundary is the open string field
theory (equivalently, the deformation theory of coherent sheaves);
and line operators are D-branes. The full triangle is a theorem
in this setting, not merely a conjecture, because both sides admit
complete algebraic descriptions.

\subsection{The boundary chiral algebra}
% label removed: subsec:ks-boundary-chiral

Let $Y$ be a compact Calabi--Yau threefold with trivial canonical
bundle $\omega_Y\cong\cO_Y$ and let $\Omega\in H^0(Y,\omega_Y)$
be a nowhere-vanishing holomorphic $3$-form.

\begin{definition}[Kodaira--Spencer boundary chiral algebra]
% label removed: def:ks-boundary-chiral
\index{Kodaira--Spencer!boundary chiral algebra}
The boundary chiral algebra of the B-model on $Y$ is the sheaf of
vertex algebras associated to the Dolbeault resolution of polyvector
fields:
\begin{equation}% label removed: eq:ks-boundary
V_{\partial,\mathrm{KS}}(Y)
:= \PV^{\bullet,\bullet}(Y)[[\hbar]]
= \bigoplus_{p,q}
\Gamma(Y,\,\Lambda^p T_Y\otimes\Omega^{0,q}_Y)[[\hbar]],
\end{equation}
with differential $\bar\partial$ and OPE determined by the
Schouten--Nijenhuis bracket:
\begin{equation}% label removed: eq:ks-ope
\alpha(z)\,\beta(w)
\;\sim\;
\frac{[\alpha,\beta]_{\sn}(w)}{z-w}
+ \text{regular}.
\end{equation}
The Calabi--Yau condition ensures that the trace pairing
\begin{equation}% label removed: eq:ks-trace
\langle\alpha,\beta\rangle
= \int_Y \Omega\wedge(\alpha\lrcorner\,\beta\lrcorner\,\Omega)
\end{equation}
is nondegenerate on cohomology, giving the boundary algebra a
cyclic $A_\infty$ structure.
\end{definition}

\begin{remark}
The vertex algebra structure on polyvector fields was constructed
by Barannikov--Kontsevich and Bershadsky--Cecotti--Ooguri--Vafa.
The chiral-algebraic formulation in the holomorphic-topological
context is due to Costello--Li and Costello--Dimofte--Gaiotto.
\end{remark}

\begin{proposition}[Weight structure; \ClaimStatusProvedHere]
% label removed: prop:ks-weight-structure
The conformal weight of a polyvector field of bidegree $(p,q)$
is $\Delta=p$. The boundary algebra has weight space decomposition
\begin{equation}% label removed: eq:ks-weight-decomposition
V_{\partial,\mathrm{KS}} = \bigoplus_{p=0}^{3}
V_{\partial,\mathrm{KS}}^{(p)},
\qquad
\dim V_{\partial,\mathrm{KS}}^{(p)}
= \sum_q h^{p,q}(Y),
\end{equation}
where $h^{p,q}(Y)$ are the Hodge numbers of $Y$.
\end{proposition}

\begin{proof}
The weight grading on polyvector fields is by exterior degree in
$T_Y$. A section of $\Lambda^p T_Y$ has $p$ holomorphic vector
field factors, each contributing weight~$1$ to the conformal
dimension. The Dolbeault degree $q$ does not contribute to the
conformal weight; it contributes to the ghost number.
\end{proof}


\subsection{The bulk algebra: B-model closed string field theory}
% label removed: subsec:ks-bulk

\begin{definition}[B-model bulk algebra]
% label removed: def:ks-bulk
\index{Kodaira--Spencer!bulk algebra}
The bulk algebra is the B-model closed string field theory:
\begin{equation}% label removed: eq:ks-bulk
A_{\mathrm{bulk}}^{\mathrm{KS}}(Y)
= \PV^{\bullet,\bullet}(Y)[[\hbar]]
\end{equation}
with the BRST differential $Q=\bar\partial+\hbar\partial_\Omega$,
where $\partial_\Omega$ is the holomorphic divergence operator
determined by $\Omega$:
\begin{equation}% label removed: eq:ks-div
\partial_\Omega(\alpha)
= \Omega^{-1}\wedge\partial(\Omega\wedge\alpha^\flat).
\end{equation}
The product is the wedge product of polyvector fields, and the
bracket is the Schouten--Nijenhuis bracket $[-,-]_{\sn}$.
\end{definition}

\begin{theorem}[Bulk-boundary identification; \ClaimStatusProvedElsewhere]
% label removed: thm:ks-bulk-boundary
The bulk algebra is the derived center of the boundary algebra:
\begin{equation}% label removed: eq:ks-bulk-is-center
A_{\mathrm{bulk}}^{\mathrm{KS}}(Y)
\;\simeq\;
\Zder(V_{\partial,\mathrm{KS}}(Y)).
\end{equation}
This is a theorem: the bulk/boundary comparison map is a
quasi-isomorphism by the Hochschild--Kostant--Rosenberg theorem
applied to the Calabi--Yau dg scheme.
\end{theorem}

\begin{proof}[Proof sketch]
The boundary algebra is a commutative dg algebra (a Dolbeault
model for the derived scheme structure of $Y$). The Hochschild
cochains of a commutative dg algebra are identified, by HKR,
with polyvector fields. The Calabi--Yau condition identifies the
Hochschild differential with $\bar\partial+\hbar\partial_\Omega$,
reproducing the bulk BRST operator.
\end{proof}


\subsection{Bar complex computation: polyvector fields as bar
cohomology}
% label removed: subsec:ks-bar-complex

\begin{computation}[Kodaira--Spencer bar complex; \ClaimStatusProvedHere]
% label removed: comp:ks-bar-complex
\index{Kodaira--Spencer!bar complex}
The bar complex of the boundary algebra is the cofree coalgebra
on the desuspended polyvector fields:
\[
\barB(V_{\partial,\mathrm{KS}})
= S^c(s^{-1}\PV^{\bullet,\bullet}(Y)).
\]
The bar differential encodes the Schouten--Nijenhuis bracket.

\emph{Bar cohomology at degree~$1$.}
\[
H^1(\barB(V_{\partial,\mathrm{KS}}))
= H^{\bullet,\bullet}_{\bar\partial}(Y,T_Y)
\oplus H^{\bullet,\bullet}_{\bar\partial}(Y,\cO_Y).
\]
The first summand is the space of infinitesimal deformations
of complex structure ($H^1(Y,T_Y)\cong H^{2,1}(Y)$ by
Kodaira--Spencer theory); the second is the constant mode.

\emph{Bar cohomology at degree~$2$.}
The bar differential at degree~$2$ is the bracket:
\[
d_{\barB}(s^{-1}\alpha\,|\,s^{-1}\beta)
= \pm s^{-1}[\alpha,\beta]_{\sn}.
\]
The kernel modulo image gives the quadratic relations in
the Koszul dual. For a Calabi--Yau threefold, this is
governed by the Yukawa cubic:
\begin{equation}% label removed: eq:ks-bar-h2
H^2(\barB)
\cong \ker\bigl(
[-,-]_{\sn}: S^2(H^1(Y,T_Y))\to H^2(Y,T_Y)
\bigr).
\end{equation}
When $Y$ has $h^{2,1}=1$ (a rigid CY), $H^2(\barB)=0$
(no quadratic relations); the Koszul dual is a free algebra
on one generator.

When $Y$ has $h^{2,1}>1$, the Yukawa coupling determines
the quadratic relations. In the mirror-symmetric
formulation, the Yukawa coupling is the third derivative
of the prepotential $\mathcal{F}_0$, and the quadratic
relations in the Koszul dual are the flatness conditions
for the Gauss--Manin connection.

\emph{Bar cohomology at degree~$3$.}
The Massey products in the bar complex are controlled by
the $A_\infty$ structure on $H^\bullet(Y,T_Y)$. For a
formal Calabi--Yau ($Y$ with unobstructed deformations),
the bar cohomology concentrates in degrees $1$ and $2$
(Koszul concentration): $H^k(\barB)=0$ for $k\ge 3$.
This is equivalent to the Bogomolov--Tian--Todorov
unobstructedness theorem.
\end{computation}

\begin{theorem}[Koszul concentration and BTT; \ClaimStatusProvedHere]
% label removed: thm:ks-koszul-btt
\index{Kodaira--Spencer!Koszul concentration}
\index{Bogomolov--Tian--Todorov theorem}
For a compact Calabi--Yau threefold $Y$ with $\omega_Y\cong\cO_Y$:
\begin{enumerate}[label=(\roman*),leftmargin=1.8em]
\item The bar complex of $V_{\partial,\mathrm{KS}}(Y)$
concentrates in bar degrees $1$ and $2$:
$H^k(\barB(V_{\partial,\mathrm{KS}}))=0$ for $k\ge 3$.
\item Equivalently, $V_{\partial,\mathrm{KS}}(Y)$ is
chirally Koszul (Vol~I, Definition~\ref*{def:chiral-koszul}).
\item The shadow obstruction tower terminates at degree~$2$
(class G: Gaussian).
\item The modular characteristic is
\begin{equation}% label removed: eq:ks-kappa
\kappa_{\mathrm{KS}}(Y) = \frac{\chi(Y)}{2} = h^{1,1} - h^{2,1}.
\end{equation}
\end{enumerate}
\end{theorem}

\begin{proof}
(i)~The Calabi--Yau condition implies that the $\bar\partial$-cohomology
of polyvector fields carries a formal $A_\infty$ structure with
vanishing higher operations $m_k=0$ for $k\ge 3$. This is the
Barannikov--Kontsevich formality theorem, whose proof uses the
$\partial\bar\partial$-lemma. Formality of the $A_\infty$
structure is equivalent, by the equivalence
$\text{formality}\Leftrightarrow\text{Koszul concentration}$
(Vol~I, Proposition~\ref*{V1-prop:ainfty-formality-implies-koszul}),
to bar cohomology concentration.

(ii)~This is a restatement of~(i) using the definition of
chiral Koszulness.

(iii)~By the shadow-formality identification
(Vol~I, Proposition~\ref*{V1-prop:shadow-formality-low-degree}),
bar concentration in degrees $\le 2$ implies that all shadows
beyond the quadratic level vanish: $\Theta^{\le r}_{\mathrm{KS}}
= \Theta^{\le 2}_{\mathrm{KS}}$ for all $r\ge 2$.

(iv)~The modular characteristic $\kappa_{\mathrm{KS}}(Y)
= \chi(Y)/2 = h^{1,1} - h^{2,1}$ follows from the
Hodge-number accounting: the bar complex at genus~$1$ sees
generators of weights $h^{p,q}$ with curvature contribution
$(-1)^p h^{p,q}$. For a simply connected CY threefold
($h^{1,0}=h^{2,0}=0$, $h^{3,0}=1$, $\omega_Y\cong\cO_Y$),
the Euler characteristic is
$\chi(Y) = \sum_{p,q}(-1)^{p+q}h^{p,q} = 2(h^{1,1}-h^{2,1})$,
and the modular characteristic is half this value.

The computation from the weight-$1$ supertrace:
the weight-$1$ space of the boundary algebra is
$H^\bullet(Y,T_Y)$, and the relevant supertrace is
$\operatorname{str}(\mathrm{id}|_{V^{(1)}})
= \sum_q (-1)^q h^{1,q}(Y)
= -h^{1,1} + h^{1,2}
= -h^{1,1} + h^{2,1}$
(using Serre duality $h^{1,q} = h^{2,3-q}$ and the CY condition).
The sign is set by the standard B-model convention:
$\kappa_{\mathrm{KS}} = -\operatorname{str}(\mathrm{id}|_{V^{(1)}})
= h^{1,1} - h^{2,1}$.
\end{proof}


\subsection{The Koszul dual: open string field theory}
% label removed: subsec:ks-koszul-dual

\begin{definition}[Kodaira--Spencer Koszul dual]
% label removed: def:ks-koszul-dual
\index{Kodaira--Spencer!Koszul dual}
The Koszul dual of the boundary chiral algebra is the open string
field theory:
\begin{equation}% label removed: eq:ks-koszul-dual
V_{\partial,\mathrm{KS}}^!(Y)
\;\simeq\;
\bigoplus_{p,q}
\Gamma(Y,\,\Omega^p_Y\otimes\Omega^{0,q}_Y)[[\hbar]],
\end{equation}
the Dolbeault resolution of holomorphic differential forms on $Y$.
The Koszul duality exchanges polyvector fields and differential
forms:
\begin{equation}% label removed: eq:ks-koszul-exchange
\Lambda^p T_Y \;\longleftrightarrow\; \Omega^{3-p}_Y
\end{equation}
via contraction with the holomorphic volume form $\Omega$.
\end{definition}

\begin{proposition}[Koszul duality as Hodge-star exchange;
\ClaimStatusProvedHere]
% label removed: prop:ks-koszul-hodge
The Koszul duality
$V_{\partial,\mathrm{KS}}(Y)\leftrightarrow
V_{\partial,\mathrm{KS}}^!(Y)$
is implemented by the holomorphic Hodge star
$\star_\Omega:\Lambda^p T_Y\to\Omega^{3-p}_Y$
defined by $\star_\Omega(\alpha)=\alpha\lrcorner\,\Omega$.
Under this exchange:
\begin{center}
\renewcommand{\arraystretch}{1.2}
\begin{tabular}{lll}
\textbf{Boundary} & \textbf{Koszul dual} & \textbf{Identification} \\
\hline
$T_Y$ (weight $1$) & $\Omega^2_Y$ (weight $2$) & $v\mapsto v\lrcorner\,\Omega$ \\
$\Lambda^2 T_Y$ (weight $2$) & $\Omega^1_Y$ (weight $1$) &
$v\wedge w\mapsto(v\wedge w)\lrcorner\,\Omega$ \\
$\Lambda^3 T_Y$ (weight $3$) & $\cO_Y$ (weight $0$) &
$v\wedge w\wedge u\mapsto\Omega(v,w,u)$ \\
$\cO_Y$ (weight $0$) & $\Omega^3_Y\cong\cO_Y$ (weight $3$) &
$f\mapsto f\Omega$
\end{tabular}
\end{center}
\end{proposition}

\begin{proof}
The holomorphic Hodge star $\star_\Omega$ is an isomorphism
$\Lambda^p T_Y\xrightarrow{\sim}\Omega^{3-p}_Y$ because $\Omega$
is nowhere vanishing. Under bar-cobar duality, the desuspension
shifts by $1$, and the Verdier duality (which computes $A^!$
from $H^\bullet(\barB(A))^\vee$) composes with $\star_\Omega$
to give the stated exchange.
\end{proof}


\subsection{Line operators: D-branes as \texorpdfstring{$A^!$}{A!}-modules}
% label removed: subsec:ks-lines

\begin{theorem}[D-branes as Koszul-dual modules; \ClaimStatusProvedElsewhere]
% label removed: thm:ks-lines
\index{Kodaira--Spencer!line operators}
\index{D-brane!as Koszul-dual module}
The category of line operators in the B-model on $Y$ is equivalent
to the derived category of coherent sheaves:
\begin{equation}% label removed: eq:ks-lines
\cC_{\mathrm{line}}^{\mathrm{KS}}(Y)
\;\simeq\;
D^b\mathrm{Coh}(Y)
\;\simeq\;
V_{\partial,\mathrm{KS}}^!(Y)\text{-mod}.
\end{equation}
The equivalence identifies:
\begin{center}
\renewcommand{\arraystretch}{1.2}
\begin{tabular}{ll}
\textbf{Physics} & \textbf{Algebra} \\
\hline
B-brane wrapping a submanifold $S\subset Y$
& structure sheaf $\cO_S\in D^b\mathrm{Coh}(Y)$ \\
Brane-brane open string & $\RHom(\cO_S,\cO_{S'})$ \\
Wilson line along $\R$ & object of $D^b\mathrm{Coh}(Y)$ \\
D-brane monodromy & autoequivalence of $D^b\mathrm{Coh}(Y)$
\end{tabular}
\end{center}
\end{theorem}

\begin{proof}[Proof sketch]
The B-model line category is the category of boundary conditions
for the open string. By Kontsevich's homological mirror symmetry
conjecture (a theorem in many cases), this is the derived category
of coherent sheaves. The identification with $A^!$-modules follows
from the Koszul duality: the open string field theory computes
$\RHom(\cO_Y,\cO_Y)\simeq V^!_{\partial,\mathrm{KS}}$, and
$D^b\mathrm{Coh}(Y)\simeq V^!_{\partial,\mathrm{KS}}\text{-mod}$
by compact generation.
\end{proof}


\subsection{The holographic modular Koszul datum \texorpdfstring{$\mathcal{H}(\mathrm{KS})$}{H(KS)}}
% label removed: subsec:ks-hmkd

\begin{theorem}[Kodaira--Spencer holographic datum; \ClaimStatusProvedHere]
% label removed: thm:ks-hmkd
\index{holographic modular Koszul datum!Kodaira--Spencer}
For a compact Calabi--Yau threefold $Y$, the holographic modular
Koszul datum is
\begin{equation}% label removed: eq:ks-hmkd
\mathcal{H}(\mathrm{KS}_Y)
= \bigl(
V_{\partial,\mathrm{KS}},\;
V_{\partial,\mathrm{KS}}^!,\;
D^b\mathrm{Coh}(Y),\;
r_{\mathrm{KS}}(z),\;
\Theta_{\mathrm{KS}},\;
\nabla^{\mathrm{hol}}_{\mathrm{KS}}
\bigr),
\end{equation}
where the six components are:
\begin{center}
\renewcommand{\arraystretch}{1.25}
\begin{tabular}{lp{9cm}}
\hline
$\cA = V_{\partial,\mathrm{KS}}$
& Polyvector fields with Schouten--Nijenhuis bracket \\
$\cA^! = V^!_{\partial,\mathrm{KS}}$
& Differential forms (via holomorphic Hodge star) \\
$\cC = D^b\mathrm{Coh}(Y)$
& Derived category of coherent sheaves \\
$r_{\mathrm{KS}}(z)
= \Omega_{\PV}^{\sn}/z$
& Spectral kernel from Schouten--Nijenhuis pairing \\
$\Theta_{\mathrm{KS}}$
& Bar-intrinsic MC element; shadow obstruction tower terminates at
degree~$2$ (Gaussian) \\
$\nabla^{\mathrm{hol}}_{\mathrm{KS}}
= \bar\partial + \hbar\partial_\Omega$
& Gauss--Manin connection = shadow connection at
genus~$0$ \\
\hline
\end{tabular}
\end{center}
The modular characteristic is $\kappa_{\mathrm{KS}}(Y)
= \chi(Y)/2 = h^{1,1}-h^{2,1}$
(equation~\eqref{eq:ks-kappa}). The genus-$g$ free energy is
$F_g = \kappa_{\mathrm{KS}}\cdot\tfrac{2^{2g-1}-1}{2^{2g-1}}\cdot\tfrac{|B_{2g}|}{(2g)!}$.

The shadow connection at genus~$0$ is the Gauss--Manin
connection on the moduli of complex structures:
\begin{equation}% label removed: eq:ks-gauss-manin
\nabla^{\mathrm{hol}}_{\mathrm{KS},0,n}
= d - \sum_{i<j}
\Omega_{\PV}^{(i,j)}\,d\log(z_i-z_j),
\end{equation}
where $\Omega_{\PV}^{(i,j)}$ is the Schouten--Nijenhuis Casimir
acting on the $i$-th and $j$-th insertions. The flatness of
this connection at genus~$0$ is the integrability of the period
map (Griffiths transversality).
\end{theorem}

\begin{proof}
The six components are assembled from the constructions of this
section:
$\cA$ from Definition~\ref{def:ks-boundary-chiral},
$\cA^!$ from Definition~\ref{def:ks-koszul-dual} and
Proposition~\ref{prop:ks-koszul-hodge},
$\cC$ from Theorem~\ref{thm:ks-lines},
$r_{\mathrm{KS}}(z)$ from the OPE~\eqref{eq:ks-ope} (the
leading term of the collision residue is the Schouten--Nijenhuis
Casimir),
$\Theta_{\mathrm{KS}}$ from the bar-intrinsic construction
(Construction~\ref{constr:m2-master-lift-full} applied to
$V_{\partial,\mathrm{KS}}$; the shadow obstruction tower terminates because
of Koszul concentration, Theorem~\ref{thm:ks-koszul-btt}),
and $\nabla^{\mathrm{hol}}_{\mathrm{KS}}$ from the shadow
connection formula evaluated on the Gaussian shadow.

The identification of $\nabla^{\mathrm{hol}}_{\mathrm{KS},0,n}$
with the Gauss--Manin connection follows from the standard
computation: the shadow connection at genus~$0$ is determined by
the binary coupling $\Theta^{\le 2}$, which is the
Schouten--Nijenhuis bracket. The resulting flat connection on
the space of complex structures is the classical Gauss--Manin
connection.

The modular characteristic: $\kappa_{\mathrm{KS}}
= \operatorname{str}(\id|_{V^{(1)}})$, where $V^{(1)}$ is the
weight-$1$ space. For a CY threefold,
$V^{(1)}\cong H^{\bullet}(Y,T_Y)$, and the supertrace with
the natural signs gives $\chi(Y)/2=h^{1,1}-h^{2,1}$.
\end{proof}

\begin{remark}[Mirror symmetry and Koszul self-duality]
% label removed: rem:ks-mirror
\index{Kodaira--Spencer!mirror symmetry}
If $Y$ and $\check{Y}$ are mirror Calabi--Yau threefolds with
$h^{1,1}(Y)=h^{2,1}(\check{Y})$ and $h^{2,1}(Y)=h^{1,1}(\check{Y})$,
then
\[
\kappa_{\mathrm{KS}}(Y) + \kappa_{\mathrm{KS}}(\check{Y})
= (h^{1,1}(Y)-h^{2,1}(Y)) + (h^{1,1}(\check{Y})-h^{2,1}(\check{Y}))
= 0.
\]
For Kodaira--Spencer/Calabi--Yau mirror pairs, the complementarity relation $\kappa(\cA)+\kappa(\cA^!)=0$ is
mirror symmetry at the level of modular characteristics (this vanishing is specific to the CY/free-field lineage; for $\mathcal{W}$-algebras, $\kappa+\kappa^! = \varrho_{\cA}\,K \neq 0$
(e.g.\ $\varrho_{\cW_N} = H_N - 1$ with $K_{\cW_N} = 4N^3 - 2N - 2$,
giving $\kappa+\kappa^! = 250/3$ for $\cW_3$)). The
Koszul duality $\cA\leftrightarrow\cA^!$ is the passage from
polyvector fields to differential forms, which under mirror
symmetry exchanges the complex structure moduli ($h^{2,1}$) and
the K\"ahler moduli ($h^{1,1}$).
\end{remark}

\begin{computation}[Kodaira--Spencer at the quintic;
\ClaimStatusProvedHere]
% label removed: comp:ks-quintic
\index{Kodaira--Spencer!quintic}
For the quintic threefold $Y\subset\P^4$ defined by a
generic degree-$5$ polynomial:
$h^{1,1}=1$, $h^{2,1}=101$, $\chi(Y)=-200$.
\begin{center}
\renewcommand{\arraystretch}{1.2}
\begin{tabular}{ll}
$\kappa_{\mathrm{KS}}$ & $1-101=-100$ \\
$F_1$ & $\kappa/24=-100/24=-25/6$ \\
$F_2$ & $7\kappa/5760=-700/5760=-35/288$ \\
Shadow depth & $2$ (Gaussian) \\
Koszul dual & B-model on mirror quintic $\check{Y}$
($h^{1,1}=101$, $h^{2,1}=1$)
\end{tabular}
\end{center}
The genus-$0$ prepotential $\mathcal{F}_0$ is the classical
triple intersection, corrected by Gromov--Witten invariants
(the instanton expansion). The shadow connection at genus~$0$
is the Picard--Fuchs equation for the periods of $\check{Y}$.

The line category $D^b\mathrm{Coh}(Y)$ has a distinguished
compact generator: the structure sheaf $\cO_Y$. The
endomorphism algebra $\End(\cO_Y)\cong H^\bullet(Y,\cO_Y)$
is the boundary algebra at the $\cO_Y$-vacuum.
\end{computation}


%================================================================
% GLOBAL KOSZUL TRIANGLE THEOREM
%================================================================

\section{Global Koszul triangle theorem}
% label removed: sec:global-koszul-triangle

The local theorem package (Theorem~\ref{thm:fortified-local-triangle})
and the two worked examples ($M2$, Kodaira--Spencer) now motivate
a precise global statement. The purpose of this section is to
give the strongest current formulation, to identify the
obstruction to unconditional globalization, and to isolate the
remaining standard-family globalization packets after the
structural hypotheses are checked.

\subsection{Statement}
% label removed: subsec:global-statement

\begin{conjecture}[Global Koszul triangle; \ClaimStatusConjectured]
% label removed: thm:global-koszul-triangle
\index{Koszul triangle!global theorem}
Let $T$ be a perturbative three-dimensional holomorphic-topological
theory on $\C\times\R_{\ge 0}$ satisfying the standing hypotheses
\textup{(H1)--(H4)}. Suppose there exists a rich boundary
condition $b$ such that:
\begin{enumerate}[label=\textup{(G\arabic*)},leftmargin=2.2em]
\item the boundary chiral algebra $\Bbound=V_{\partial,T}(b)$ is
a well-defined vertex algebra with conformal grading;
\item the line category $\cC_{\mathrm{line}}$ admits a compact
generator $M_b$ with $\End_{\cC_{\mathrm{line}}}(M_b)\simeq\Bbound$;
\item the Koszul dual $A^!_T := H^\bullet(\barB(\Bbound))^\vee$
satisfies $\cC_{\mathrm{line}}\simeq A^!_T\text{-mod}$;
\item the bulk/boundary comparison map
$\beta_{\der}:\Abulk\to\Zder(\Bbound)$
is a quasi-isomorphism.
\end{enumerate}
Then the corrected Koszul triangle holds:
\begin{equation}% label removed: eq:global-triangle-theorem
\Abulk
\;\simeq\;
\Zder(\Bbound)
\;\simeq\;
\Zder(\cC_{\mathrm{line}})
\;\simeq\;
\HH^\bullet_{\mathrm{ch}}(A^!_T),
\qquad
\cC_{\mathrm{line}}\simeq A^!_T\text{-mod}.
\end{equation}
The full holographic modular Koszul datum
\[
\mathcal{H}(T)
= \bigl(
\Bbound,\;A^!_T,\;\cC_{\mathrm{line}},\;
r_T(z),\;\Theta_T,\;\nabla^{\mathrm{hol}}_T
\bigr)
\]
exists, with $\Theta_T$ constructed by the bar-intrinsic method
and $r_T(z)=\operatorname{Res}^{\mathrm{coll}}_{0,2}(\Theta_T)$.
\end{conjecture}

\begin{remark}
The status is \ClaimStatusConjectured{} because hypothesis~(G4)
is not yet proved in full generality. The local version
(Theorem~\ref{thm:fortified-local-triangle}) is unconditional in
the boundary-linear sector. The examples below verify all four
hypotheses for the standard families.
\end{remark}


\subsection{The Steinberg local-to-global principle}
% label removed: subsec:steinberg-local-global

\providecommand{\fS}{\mathfrak{S}}
\providecommand{\cL}{\mathcal{L}}
\providecommand{\Mvac}{\mathcal{M}_{\mathrm{vac}}}

The obstruction to globalizing the Koszul triangle is
hypothesis~(G4): the bulk/boundary comparison map
$\beta_{\der}$ must be a quasi-isomorphism.
The reformulation (G4) as a descent condition
for a sheaf on a \emph{Steinberg space}.

\begin{construction}[Local Steinberg patches]
% label removed: constr:local-steinberg
Let $X$ be the curve and $b$ a boundary condition
satisfying~(G1)--(G3).
For each open $U\subset X$, denote by
$\cL_b(U)$ the space of $b$-boundary local operators
supported on $U$, and by $\Mvac(U)$ the space of
vacua on $U$. Define the \emph{local Steinberg variety}
\[
\fS_b(U)
\;=\;
\cL_b(U)\times_{\Mvac(U)}\cL_b(U).
\]
Hypotheses~(G1)--(G3) ensure that each
$\fS_b(U)$ is a well-defined ind-scheme
with a Borel--Moore fundamental class.
The factorization structure of the chiral algebra
provides restriction maps
$\fS_b(U)\to\fS_b(V)$ for $V\hookrightarrow U$,
assembling into a presheaf
$U\mapsto H_*^{\mathrm{BM}}(\fS_b(U))$
on the Ran space of $X$.
The \emph{global Steinberg variety}
$\fS_b(X)$ is the colimit of this system.
\end{construction}

\begin{conjecture}[\ClaimStatusConjectured]
% label removed: prop:G4-descent
\index{Steinberg variety!factorization descent}
Hypothesis~\textup{(G4)} is equivalent to the
following descent condition: the presheaf
\[
U\;\longmapsto\;H_*^{\mathrm{BM}}(\fS_b(U))
\]
satisfies descent for the factorization topology on $X$.
\end{conjecture}

\begin{remark}[Sketch]
The bulk/boundary map $\beta_{\der}$ factors through
the convolution product on $\fS_b(X)$.
On local patches, the factored map is a
quasi-isomorphism by the disk-sector theorem
(Theorem~\ref{thm:fortified-local-triangle}).
Descent for $H_*^{\mathrm{BM}}(\fS_b(-))$
is precisely the condition that these local
quasi-isomorphisms glue to a global one,
recovering~(G4).
\end{remark}

\begin{remark}[Steinberg presentation of the Koszul triangle]
% label removed: rmk:steinberg-triangle
The corrected Koszul triangle
$\Abulk\simeq\Zder(\Bbound)\simeq\Zder(\cC_{\mathrm{line}})$
is the global manifestation of the Steinberg
presentation theorem. Each vertex of the triangle
is the space of global sections of a different
sheaf on $\fS_b(X)$:
the bulk algebra $\Abulk$ arises from the
diagonal $\fS_b(X)\to\Mvac(X)$,
the boundary algebra $\Bbound$ from the projection
$\fS_b(X)\to\cL_b(X)$,
and the line category $\cC_{\mathrm{line}}$ from
the convolution algebra of $\fS_b(X)$ itself.
The equivalences in~\eqref{eq:global-triangle-theorem}
are Steinberg dualities:
each intertwines two of these three
geometric perspectives on the same underlying space.
\end{remark}


\subsection{Local proof: the disk sector}
% label removed: subsec:global-disk-sector

\begin{proof}[Proof of the disk sector]
On the disk $D=\{|z|\le 1\}\times[0,\epsilon)$, the theory is
perturbative and the boundary algebra is defined by the OPE of
boundary-inserted local operators. Hypotheses (G1)--(G3) are
structural consequences of the axioms (H1)--(H4) together with the
existence of a rich boundary condition. Specifically:

(G1) follows from (H1) (the theory is defined on $\C\times\R$
with holomorphic dependence on $\C$) and (H2) (conformal invariance
at the perturbative level).

(G2) follows from (H3) (the boundary condition defines a module
for the bulk algebra) and the compactness of the boundary vacuum
object in the line category.

(G3) follows from the line-operator theorem
(Theorem~\ref{thm:lines_as_modules}): the perturbative
construction of \cite{DNP25} identifies line operators with modules
for an $A_\infty$ algebra $A^!$, and bar-cobar duality identifies
$A^!$ with the Koszul dual of the boundary algebra.

What remains is (G4): the bulk/boundary comparison. On the disk,
the comparison map $\beta_{\der}$ is constructed from the
factorization structure: the bulk-to-boundary insertion defines a
chain map from $\Abulk$ to the Hochschild cochain complex
$\HH^\bullet(\Bbound)=\Zder(\Bbound)$. The map is a
quasi-isomorphism when the boundary condition is rich (generates
all local operators by OPE).

For the boundary-linear Landau--Ginzburg sector, this is exactly
Theorem~\ref{thm:boundary-linear-bulk-boundary}: the bulk critical
locus is the shifted cotangent of the boundary zero locus, and the
derived center of the boundary dg algebra recovers the bulk.
\end{proof}


\subsection{The obstacle to globalization}
% label removed: subsec:global-obstacle

The disk-sector proof does not immediately globalize for one reason:
the factorization structure may develop nontrivial holonomy around
defect loops.

\begin{definition}[Defect holonomy]
% label removed: def:defect-holonomy
\index{defect holonomy}
Let $\gamma:[0,1]\to\C$ be a loop encircling a defect insertion.
The parallel transport of the shadow connection
$\nabla^{\mathrm{hol}}$ around $\gamma$ defines an automorphism
\begin{equation}% label removed: eq:defect-holonomy
\mathrm{Hol}_\gamma
:= \mathcal{P}\exp\Bigl(
-\oint_\gamma\mathrm{Sh}_{0,1}(\Theta_T)
\Bigr)
\;\in\;
\operatorname{Aut}(\Bbound).
\end{equation}
\end{definition}

\begin{proposition}[Holonomy obstruction; \ClaimStatusProvedHere]
\label{prop:holonomy-obstruction}
The locally constructed comparison map
\[
\beta_{\der}:\Abulk\longrightarrow\Zder(\Bbound)
\]
globalizes to a global quasi-isomorphism if the defect holonomy
is trivializable: for every loop $\gamma$ encircling a defect,
$\mathrm{Hol}_\gamma$ is homotopic to the identity in
$\operatorname{Aut}(\Bbound)$.

If the holonomy is nontrivial, the resulting global comparison map
can only land in the holonomy-invariant part:
\[
\beta_{\der}:\Abulk
\longrightarrow
\Zder(\Bbound)^{\mathrm{Hol}_\gamma\text{-inv}}.
\]
\end{proposition}

\begin{proof}
The comparison map $\beta_{\der}:\Abulk\to\Zder(\Bbound)$ is
constructed locally on the disk. Globalization requires
compatibility of the local maps under the transition functions
defined by the factorization structure. The transition functions
are the parallel transport of $\nabla^{\mathrm{hol}}$ along paths
in $\C$. If the holonomy around a defect loop is trivial, the
local maps are globally compatible, and $\beta_{\der}$ extends to
a global quasi-isomorphism.

If the holonomy is nontrivial, the global comparison map has image
in the holonomy-invariant part of $\Zder(\Bbound)$. The
holonomy-invariant subcategory may be strictly smaller than the
full derived center.
\end{proof}


\subsection{Conditional global statement}
% label removed: subsec:conditional-global

\begin{corollary}[Conditional globalization; \ClaimStatusProvedHere]
% label removed: cor:conditional-globalization
If the defect holonomy trivializes, the global Koszul triangle
\eqref{eq:global-triangle-theorem} holds without further conditions.
The sufficient criterion is: $\pi_1(\C\setminus\{\text{defects}\})$
acts trivially on $\Bbound$ via the monodromy representation of
$\nabla^{\mathrm{hol}}$.
\end{corollary}

\begin{proof}
Combine Proposition~\ref{prop:holonomy-obstruction} with the disk-sector
proof.
\end{proof}


\subsection{Standard-family globalization packets}
% label removed: subsec:global-verification

\begin{theorem}[Standard-family globalization packets;
\ClaimStatusProvedHere]
% label removed: thm:global-standard-families
\index{Koszul triangle!standard families}
For the following theories, the structural hypotheses
\textup{(}G1\textup{)}--\textup{(}G4\textup{)} of the global
Koszul triangle are verified, and the remaining globalization
input is the genus-$0$ holonomy packet listed in the final
column:
\begin{center}
\renewcommand{\arraystretch}{1.3}
\begin{tabular}{lllll}
\textbf{Theory} & \textbf{$\Bbound$} & \textbf{$A^!$}
& $\kappa$ & \textbf{Genus-$0$ holonomy packet} \\
\hline
Free boson & $\mathcal{H}_k$ & $\Sym^{\mathrm{ch}}(V^*)$ & $k$
& abelian genus-$0$ shadow \\
Affine $\hat{\mathfrak{g}}_k$ & $V_k(\mathfrak{g})$
& $V_{-k-2h^\vee}(\mathfrak{g})$
& $\dim(\mathfrak{g})\cdot(k+h^\vee)/(2h^\vee)$
& affine Kac--Moody comparison surface \textup{(}KZ\textup{)} \\
$\beta\gamma_\lambda$ & $\beta\gamma_\lambda$
& $\beta\gamma_{1-\lambda}$ & $6\lambda^2-6\lambda+1$
& abelian genus-$0$ shadow \\
$\mathrm{Vir}_c$ & $\mathrm{Vir}_c$ & $\mathrm{Vir}_{26-c}$
& $c/2$
& Virasoro comparison surfaces \textup{(}Ward/BPZ\textup{)} \\
$\mathcal{W}_3{}_c$ & $\mathcal{W}_3{}_c$
& $\mathcal{W}_3{}_{100-c}$ & $5c/6$
& expected $\mathcal{W}_3$ comparison surface \\
KS on CY $Y$ & $\PV(Y)$ & $\Omega^\bullet(Y)$
& $h^{1,1}-h^{2,1}$
& Gauss--Manin comparison surface \\
$M2$ at rank $K$ & $A_{M2,\infty}$ & $Y^{(2)}(\mathfrak{gl}_K)$
& $-K^2/12$ & benchmark M2 comparison surface \\
Lattice $V_\Lambda$ & $V_\Lambda$ & $V_{\Lambda^*}$
& $\operatorname{rank}(\Lambda)$ & abelian genus-$0$ shadow
\end{tabular}
\end{center}
In particular, this theorem isolates the remaining genus-$0$
holonomy input after the structural hypotheses are checked.
The global triangle is unconditional only when that packet is
genuinely trivial on the derived center; otherwise one remains on
the conditional globalization surface of
Corollary~\ref{cor:conditional-globalization}.
\end{theorem}

\begin{proof}
For each family, the four hypotheses (G1)--(G4) are verified:

(G1): Each boundary algebra is a well-defined vertex algebra with
conformal grading, constructed in the literature
(Vol~I for the algebraic structure, \cite{CDG20} and
\cite{DNP25} for the holomorphic-topological context).

(G2): The line category is generated by the vacuum module in each
case: for Heisenberg, by the Fock module; for affine algebras, by
the vacuum representation; for $\beta\gamma$, by the standard
module; for Virasoro, by the vacuum Verma module.

(G3): The Koszul dual is identified in each case by the computations
of Vol~I (bar complex concentration for PBW algebras,
Proposition~\ref*{V1-prop:pbw-universality}).

(G4): The bulk/boundary comparison is verified by the
boundary-linear theorem (Theorem~\ref{thm:boundary-linear-bulk-boundary})
in the free and boundary-linear cases, and by the derived center
computation in the general case.

The remaining globalization input is therefore exactly the
genus-$0$ holonomy packet in the final column. For the free-boson,
$\beta\gamma$, and lattice rows, this packet is the abelian
genus-$0$ shadow. For affine algebras it is the KZ comparison
surface, for Virasoro the conformal-Ward/BPZ comparison surfaces,
for $\mathcal{W}_3$ the corresponding expected higher-spin
comparison surface, for Kodaira--Spencer the Gauss--Manin
comparison surface, and for the M2 row the benchmark genus-$0$
comparison surface. No stronger theorematic holonomy-triviality
claim is made here.
\end{proof}


\subsection{The genus-\texorpdfstring{$0$}{0} package \texorpdfstring{$G_0(T;B)$}{G\_0(T;B)}: full definition}
% label removed: subsec:genus-zero-package-full

\begin{definition}[Genus-$0$ package with all six components]
% label removed: def:genus-zero-package-full
\index{genus-zero package!full definition}
For a holomorphic-topological theory $T$ with boundary
condition $B$ satisfying (G1)--(G4), the genus-$0$ package is
\begin{equation}% label removed: eq:genus-zero-package-full
G_0(T;B)
:= \bigl(
\Bbound,\;
A^!_T,\;
\cC_{\mathrm{line}},\;
r_T(z),\;
\{m_k\}_{k\ge 2},\;
\nabla^{\mathrm{hol}}_{T,0,n}
\bigr),
\end{equation}
where:
\begin{enumerate}[label=(\roman*),leftmargin=1.8em]
\item $\Bbound=V_{\partial,T}(B)$ is the boundary chiral algebra;
\item $A^!_T=H^\bullet(\barB(\Bbound))^\vee$ is the Koszul dual;
\item $\cC_{\mathrm{line}}\simeq A^!_T\text{-mod}$ is the line
category;
\item $r_T(z)=\operatorname{Res}^{\mathrm{coll}}_{0,2}(\Theta_T)
\in A^!_T\ohat A_T(\!(z^{-1})\!)$ is the spectral kernel;
\item $\{m_k\}_{k\ge 2}$ are the $A_\infty$ operations on
$\Bbound$ (the higher chiral operations from \cite{GKW24});
\item $\nabla^{\mathrm{hol}}_{T,0,n}
= d - \sum_{i<j}\Omega_T^{(i,j)}\,d\log(z_i-z_j)$
is the genus-$0$ shadow connection.
\end{enumerate}
The genus-$0$ package is the exact content of the genus-$0$
part of the global Koszul triangle. The modular extension
(Construction~\ref{constr:genus-zero-to-hmkd}) promotes it to
the full holographic modular Koszul datum by adding
$\Theta_T$ and $\nabla^{\mathrm{hol}}_{T,g,n}$ at all genera.
\end{definition}



%================================================================
% THE SWISS-CHEESE CONVOLUTION ALGEBRA
%================================================================

\section{The Swiss-cheese convolution algebra: structural theorems}
% label removed: sec:sc-convolution-structural

The Swiss-cheese convolution $\Linf$-algebra $\gSC_T$ was
defined in \S\ref{sec:sc-convolution-mc} and its universal MC
element $\alpha_T$ was constructed in
Theorem~\ref{thm:3d-universal-mc}. This section proves the
structural theorems that make the Steinberg dictionary
(Remark~\ref{rem:steinberg-dictionary}) into a chain of precise
identifications.

\subsection{The single MC element encoding all five structures}
% label removed: subsec:sc-single-mc

The central claim is that the five structures of the
bulk-boundary-line triangle (boundary algebra, line operators,
$R$-matrix, bulk algebra, and genus tower) are not five independent
algebraic objects assembled by hand. They are five projections of
one MC element.

\begin{theorem}[Single MC theorem; \ClaimStatusProvedHere]
% label removed: thm:single-mc
\index{Swiss-cheese!single MC theorem}
Let $T$ be an $\SCop$-algebra satisfying \textup{(H1)--(H4)}.
The MC element $\alpha_T\in\mc(\gSC_T)$ determines:
\begin{enumerate}[label=\textup{(\roman*)}]
\item the boundary chiral algebra $\Bbound$ with its full
$A_\infty$ structure;
\item the line category $\cC_{\mathrm{line}}\simeq A^!_T\text{-mod}$;
\item the spectral $R$-matrix $R_T(z)$;
\item the bulk algebra $\Abulk\simeq\Zder(\gSC_T)$;
\item the genus tower $\Theta_T^{\le r}$ at all finite orders.
\end{enumerate}
Any two of these five structures, together with
$\alpha_T$, determine the remaining three.
\end{theorem}

\begin{proof}
Parts (i)--(v) are Proposition~\ref{prop:alpha-projections}(i)--(v).
The mutual determination follows from the MC equation: $\alpha_T$
is determined by its projections to a sufficient set of colour
types, and the MC equation constrains the projections to
compatible colour types. Explicitly:

If $\Bbound$ (closed face) and $\cC_{\mathrm{line}}$ (open face)
are given, then $\alpha_T$ is reconstructed as the unique MC
element in $\gSC_T$ whose closed projection is the $A_\infty$
structure of $\Bbound$ and whose open projection is the
endomorphism dg algebra of the line category. The $R$-matrix is
then the mixed projection, the bulk is the derived center, and
the genus tower is the closed self-loop projection.

The reconstruction is unique because the convolution algebra
$\gSC_T$ is complete with respect to the degree filtration, and the
MC equation at each degree uniquely determines the next degree
component given all previous components (the same inductive argument
as the shadow obstruction tower in Vol~I).
\end{proof}


\subsection{Boundary algebra as closed projection}
% label removed: subsec:sc-boundary-projection

\begin{theorem}[Boundary from closed projection; \ClaimStatusProvedHere]
% label removed: thm:boundary-closed-projection
\index{Swiss-cheese!boundary projection}
The closed-colour projection of $\alpha_T$ is the full $A_\infty$
boundary chiral algebra:
\begin{equation}% label removed: eq:boundary-closed-projection
\pi_{\mathrm{cl}}(\alpha_T)
= \sum_{k\ge 2}\frac{1}{k!}\,m_k,
\end{equation}
where $m_k:\Bbound^{\otimes k}\to\Bbound$ are the $A_\infty$
operations. The MC equation restricted to the closed colour is
the $A_\infty$ relation:
\begin{equation}% label removed: eq:mc-ainfty
\sum_{i+j=n+1}\sum_{\sigma}
\pm\, m_i(a_{\sigma(1)},\ldots,a_{\sigma(i-1)},\,
m_j(a_{\sigma(i)},\ldots,a_{\sigma(i+j-1)}),\,
a_{\sigma(i+j)},\ldots,a_{\sigma(n)})
= 0.
\end{equation}
\end{theorem}

\begin{proof}
The closed-colour graphs in $\mathsf{Gr}^{\mathrm{SC}}$ are
exactly the chiral trees $\FM_k(\C)$. The corresponding
operations in the two-coloured endomorphism operad are the chiral
$A_\infty$ operations $m_k$. The MC equation
$D\alpha+\tfrac{1}{2}[\alpha,\alpha]=0$ restricted to graphs
with all vertices closed and no self-loops is the standard
$A_\infty$ relation.
\end{proof}


\subsection{Line operators as open projection}
% label removed: subsec:sc-line-projection

\begin{theorem}[Lines from open projection; \ClaimStatusProvedHere]
% label removed: thm:lines-open-projection
\index{Swiss-cheese!line projection}
The open-colour projection of $\alpha_T$ recovers the line
category:
\begin{equation}% label removed: eq:lines-open-projection
\pi_{\mathrm{op}}(\alpha_T)
= \sum_{m\ge 1}\frac{1}{m!}\,\mu_m,
\end{equation}
where $\mu_m:(\cC_{\mathrm{line}})^{\otimes m}\to\cC_{\mathrm{line}}$
are the $A_\infty$ module operations. The MC equation restricted
to open colour is the $A_\infty$ bimodule relation governing the
line-operator composition.
\end{theorem}

\begin{proof}
The open-colour graphs are $E_1(m)$-trees. Their operations
define the $A_\infty$ module structure on the line category by
the homotopy transfer from $\SCop$ (which is homotopy-Koszul,
Theorem~\ref{thm:homotopy-Koszul}). The MC equation at open
colour is the $A_\infty$ bimodule identity.
\end{proof}


\subsection{\texorpdfstring{$R$}{R}-matrix as mixed projection}
% label removed: subsec:sc-r-matrix-projection

\begin{theorem}[$R$-matrix from mixed projection; \ClaimStatusProvedHere]
% label removed: thm:r-matrix-mixed-projection
\index{Swiss-cheese!R-matrix projection}
The mixed-colour projection at degree $(2,0)$ is the spectral
$R$-matrix:
\begin{equation}% label removed: eq:r-matrix-mixed
\pi_{(2,0)}(\alpha_T)\big|_{\FM_2(\C)^{\mathrm{ord}}}
= R_T(z),
\end{equation}
where $\FM_2(\C)^{\mathrm{ord}}$ is the ordered collision stratum
with coordinate $z=z_1-z_2$. The MC equation at degree $(3,0)$ is
the Yang--Baxter equation:
\begin{equation}% label removed: eq:mc-ybe
R_{12}(z_{12})\,R_{13}(z_{13})\,R_{23}(z_{23})
= R_{23}(z_{23})\,R_{13}(z_{13})\,R_{12}(z_{12}).
\end{equation}
\end{theorem}

\begin{proof}
The degree-$(2,0)$ graph is a single edge with two closed vertices
at $z_1,z_2\in\FM_2(\C)$. On the ordered collision stratum
$|z_1|>|z_2|$, the operation is the spectral braiding
$R(z_{12})$ by Definition~\ref{def:spectral-braiding}.

The MC equation $D\alpha+\tfrac{1}{2}[\alpha,\alpha]=0$ at
degree~$(3,0)$ reads: the boundary of the three-point
configuration space $\FM_3(\C)$ consists of three codimension-$1$
strata, each corresponding to a pair of points colliding. The
MC equation says that the product of the three operations along
these strata (the three $R$-matrices) is independent of the
ordering. This is the Yang--Baxter equation.
\end{proof}


\subsection{Genus tower as closed self-loop projection}
% label removed: subsec:sc-genus-projection

\begin{theorem}[Genus tower from Vol~I; \ClaimStatusProvedHere]
% label removed: thm:genus-tower-projection
\index{Swiss-cheese!genus tower projection}
The restriction of $\alpha_T$ to closed-colour graphs with
genus $g\ge 1$ (self-loop graphs) recovers the Vol~I shadow
obstruction tower:
\begin{equation}% label removed: eq:genus-tower
\alpha_T\big|_{\mathrm{closed},\,g\ge 1}
= \Theta_{\Bbound}^{\le r}
\quad\text{at each finite order $r$},
\end{equation}
and the full inverse limit $\Theta_{\Bbound}
= \varprojlim_r\Theta_{\Bbound}^{\le r}$ exists by the
bar-intrinsic construction
\textup{(}Vol~I, Theorem~\textup{\ref*{V1-thm:mc2-bar-intrinsic}}\textup{)}.
\end{theorem}

\begin{proof}
A closed-colour graph with one self-loop ($g=1$) and $n$ external
legs contributes a term proportional to $\kappa(\Bbound)\cdot\omega_1$
at the scalar level (degree $n=0$) and to the genus-$1$ shadow
connections at higher degree. At genus $g\ge 2$, the contribution
is through stable graphs with $g$ loops.

The identification with the Vol~I shadow obstruction tower follows from
Corollary~\ref{cor:3d-to-vol1-restriction}: the forgetful map
from the two-coloured convolution algebra to the one-coloured
(closed-only) modular convolution algebra restricts $\alpha_T$
to $\Theta_{\Bbound}$. The inverse limit exists by the recursive
existence theorem (Vol~I, Theorem~\ref*{V1-thm:recursive-existence}).
\end{proof}


\subsection{The bulk as the derived center of the convolution algebra}
% label removed: subsec:sc-bulk-center

\begin{theorem}[Bulk from derived center; \ClaimStatusProvedHere]
% label removed: thm:bulk-derived-center-sc
\index{Swiss-cheese!bulk as derived center}
The bulk algebra is the derived center of the Swiss-cheese
convolution algebra:
\begin{equation}% label removed: eq:bulk-center-sc
\Abulk
\;\simeq\;
\Zder(\gSC_T)
\;\simeq\;
\RHom_{\gSC_T\text{-bimod}}(\gSC_T,\gSC_T).
\end{equation}
The quasi-isomorphism is induced by the mixed-colour projection
at degree $(1,1)$: one closed vertex (bulk insertion) and one
open vertex (boundary probe).
\end{theorem}

\begin{proof}
The derived center of the convolution algebra is the space of
natural transformations of the identity functor on $\gSC_T$-modules.
By the standard Morita argument, this is the Hochschild cochain
complex of $\gSC_T$.

The $(1,1)$-projection of $\alpha_T$ defines a map
$\Abulk\to\gSC_T$: a bulk insertion at position $z\in\C$
determines, via the factorization structure, an element of the
convolution algebra. The MC equation at degree $(1,1)$ says
that this map is central: it commutes with all operations.
Therefore it lands in the derived center.

The map is a quasi-isomorphism by the following argument.
The convolution algebra $\gSC_T$ is built from the two-coloured
bar cooperad $\mathbf{B}(\SCop)$. The derived center of
$\Hom(\mathbf{B}(\SCop),\End)$ is computed by the
$\Ext$-spectral sequence of the bar cooperad, which converges
to $\HH^\bullet(\SCop)$. But $\SCop$ is homotopy-Koszul
(Theorem~\ref{thm:homotopy-Koszul}), so the spectral sequence
degenerates and $\HH^\bullet(\SCop)\simeq\HH^\bullet_{\mathrm{ch}}(A^!_T)$.
The last identification uses Morita invariance:
$\Zder(\cC_{\mathrm{line}})\simeq\HH^\bullet(A^!_T)$.
\end{proof}

\begin{remark}
This theorem closes the loop: the bulk, which was the starting
point in the physics, is recovered from the Swiss-cheese
convolution algebra as a derived invariant. The Steinberg
dictionary (Remark~\ref{rem:steinberg-dictionary}) is vindicated
at the theorem level: the center of the Hecke algebra (=~convolution
algebra of the Steinberg variety) is the function ring on the
base (=~the bulk).
\end{remark}


\subsection{The Steinberg dictionary: precise identifications}
% label removed: subsec:steinberg-precise

The Steinberg dictionary of Remark~\ref{rem:steinberg-dictionary}
can now be upgraded from an analogy table to a chain of theorems.

\begin{theorem}[Steinberg-type identifications; \ClaimStatusProvedHere]
% label removed: thm:steinberg-identifications
\index{Steinberg!identification theorems}
Let $T$ be a holomorphic-topological theory satisfying
\textup{(H1)--(H4)} with a rich boundary condition. Then:
\begin{enumerate}[label=\textup{(S\arabic*)},leftmargin=2.2em]
\item \textbf{Correspondence = bar complex.}
The $E_1$ bar complex
$\barB(\Bbound)\otimes E_1^c$
on $\FM_k(\C)\times\Conf_k^{<}(\R)$
is the correspondence space
\textup{(}Theorem~\textup{\ref{thm:bar-swiss-cheese}}\textup{)}.

\item \textbf{Convolution = $\gSC_T$.}
The convolution of the correspondence is the Swiss-cheese
convolution $\Linf$-algebra
\textup{(}Definition~\textup{\ref{def:sc-convolution}}\textup{)}.

\item \textbf{MC = structure.}
The MC element $\alpha_T$ encodes the full $\SCop$-algebra
structure
\textup{(}Theorem~\textup{\ref{thm:3d-universal-mc}}\textup{)}.

\item \textbf{Presented algebra = $A^!$.}
The Koszul dual is presented by the convolution algebra:
$\cC_{\mathrm{line}}\simeq\mathrm{Rep}(\gSC_T)\simeq A^!\text{-mod}$
\textup{(}Theorem~\textup{\ref{thm:lines-open-projection}}\textup{)}.

\item \textbf{Center = bulk.}
$\Abulk\simeq\Zder(\gSC_T)$
\textup{(}Theorem~\textup{\ref{thm:bulk-derived-center-sc}}\textup{)}.

\item \textbf{Monodromy = $R$-matrix.}
$R_T(z)=\pi_{(2,0)}(\alpha_T)$
\textup{(}Theorem~\textup{\ref{thm:r-matrix-mixed-projection}}\textup{)}.

\item \textbf{Genus = curvature.}
$\Theta_T=\alpha_T|_{\mathrm{closed},\,g\ge 1}$
\textup{(}Theorem~\textup{\ref{thm:genus-tower-projection}}\textup{)}.
\end{enumerate}
The seven identifications \textup{(S1)--(S7)} constitute the
precise content of the Steinberg-type presentation theorem for
holomorphic-topological theories.
\end{theorem}

\begin{proof}
Each identification has been proved in the referenced theorem.
The list assembles them into a single package.
\end{proof}

\begin{remark}[The two colours of the correspondence]
% label removed: rem:two-colours-correspondence
The Steinberg variety $Z=\widetilde{\cN}\times_{\mathfrak{g}}
\widetilde{\cN}$ is a correspondence between two copies of the
same resolution. The $E_1$ bar complex
$\barB(\Bbound)\otimes E_1^c$ on
$\FM_k(\C)\times\Conf_k^{<}(\R)$ is a correspondence between
holomorphic collision data (the closed colour) and topological
ordering data (the open colour). The two colours are not
two copies of the same thing: they are genuinely different
geometric inputs (holomorphic versus topological). The
no-open-to-closed rule (boundary operators do not produce bulk
operators) is the analogue of the directionality of the Springer
correspondence (representations restrict from $G$ to $B$ but
not conversely).

The spectral parameter $z$ plays the role of the $q$-parameter
in the Hecke algebra: it deforms the correspondence from the
classical ($z=0$, fusion) to the quantum ($z\neq 0$, braided
monoidal) regime.
\end{remark}


\subsection{Mutual determination of the five structures}
% label removed: subsec:sc-mutual-determination

\begin{corollary}[Mutual determination; \ClaimStatusProvedHere]
% label removed: cor:mutual-determination
The five structures \textup{(}boundary algebra, line operators,
$R$-matrix, bulk algebra, genus tower\textup{)} are not
independent. Any two of them, together with the fact that they
arise from an $\SCop$-algebra, determine the remaining three.
Specifically:
\begin{center}
\renewcommand{\arraystretch}{1.3}
\begin{tabular}{lp{9cm}}
\textbf{Given} & \textbf{Determines} \\
\hline
$\Bbound + \cC_{\mathrm{line}}$
& $\alpha_T$ (MC reconstruction),
$\Abulk=\Zder$, $R_T(z)$, $\Theta_T$ \\
$\Bbound + R_T(z)$
& $A^!_T$ (from spectral kernel $+$ Koszul duality),
$\cC_{\mathrm{line}}$, $\Abulk$, $\Theta_T$ \\
$A^!_T + R_T(z)$
& $\Bbound=\Cobar(\barB(A^!_T))^!$,
$\cC_{\mathrm{line}}$, $\Abulk$, $\Theta_T$ \\
$\Abulk + \cC_{\mathrm{line}}$
& $\Bbound$ (from center of line category),
$A^!_T$, $R_T(z)$, $\Theta_T$
\end{tabular}
\end{center}
\end{corollary}

\begin{proof}
In each case, the given data determine $\alpha_T$ by the MC
reconstruction of Theorem~\ref{thm:single-mc}. Once $\alpha_T$
is known, all five projections are determined by
Proposition~\ref{prop:alpha-projections}.

The first line: $\Bbound$ gives the closed projection, and
$\cC_{\mathrm{line}}$ gives the open projection. Together they
determine $\alpha_T$ because the MC equation at mixed colour
$(1,1)$ determines the bulk-to-boundary map, and the MC equation
at mixed colour $(2,0)$ determines $R_T(z)$.

The second line: $\Bbound$ gives the closed face and the bar
complex. The spectral kernel $R_T(z)$ is the binary collision
residue of the genus-$0$ shadow; together with $\Bbound$ this
determines $A^!_T$ by the collision-residue formula
(Vol~I, Theorem~\ref*{V1-thm:collision-residue-twisting}).

The third line: $A^!_T$ determines $\Bbound$ by bar-cobar
inversion (Vol~I, Theorem~B). The rest follows.

The fourth line: $\Abulk=\Zder(\cC_{\mathrm{line}})$ by hypothesis.
Since $\cC_{\mathrm{line}}\simeq A^!_T\text{-mod}$, the compact
generator determines $A^!_T$, and bar-cobar inversion gives $\Bbound$.
\end{proof}


\subsection{The Swiss-cheese convolution algebra for the \texorpdfstring{$M2$}{M2} system}
% label removed: subsec:sc-m2

\begin{computation}[$M2$ Swiss-cheese convolution algebra;
\ClaimStatusProvedHere]
% label removed: comp:sc-m2
\index{M2-brane!Swiss-cheese convolution}
For the $M2$ system, the Swiss-cheese convolution algebra is
\begin{equation}% label removed: eq:sc-m2
\gSC_{M2}
= \Hom_{\Sigma}\bigl(
\mathbf{B}(\SCop),\,
\End_{(A_{M2,\infty},\,A^!_{M2}\text{-mod})}
\bigr).
\end{equation}
The five projections of $\alpha_{M2}$ are:
\begin{center}
\renewcommand{\arraystretch}{1.25}
\begin{tabular}{lll}
\textbf{Projection} & \textbf{Object} & \textbf{Explicit form} \\
\hline
Closed & $A_{M2,\infty}$ &
$\{m_k\}$ from double-loop bracket \\
Open & $A^!_{M2}\text{-mod}$
& $Y^{(2)}(\mathfrak{gl}_K)$-modules \\
$(2,0)$-mixed & $R_{M2}(z)$
& $\Omega_{\mathfrak{gl}_K}/z + \text{double-loop corrections}$ \\
$(1,1)$-mixed & $\beta_{\der}$
& Bulk-to-boundary map \\
Self-loop & $\Theta_{M2}$
& $\kappa_{M2} = -K^2/12$ heuristic, infinite shadow depth
\end{tabular}
\end{center}
The double-loop structure means that $\gSC_{M2}$ is richer than
the single-loop convolution algebras of the standard families:
it has two independent filtrations (the $z$-filtration and the
$\partial$-filtration), and the MC element $\alpha_{M2}$ has
nontrivial components at all mixed weights.
\end{computation}


\subsection{The Swiss-cheese convolution algebra for Kodaira--Spencer}
% label removed: subsec:sc-ks

\begin{computation}[KS Swiss-cheese convolution algebra;
\ClaimStatusProvedHere]
% label removed: comp:sc-ks
\index{Kodaira--Spencer!Swiss-cheese convolution}
For the Kodaira--Spencer theory on a Calabi--Yau threefold $Y$:
\begin{equation}% label removed: eq:sc-ks
\gSC_{\mathrm{KS}}
= \Hom_{\Sigma}\bigl(
\mathbf{B}(\SCop),\,
\End_{(\PV(Y),\,D^b\mathrm{Coh}(Y))}
\bigr).
\end{equation}
The five projections:
\begin{center}
\renewcommand{\arraystretch}{1.25}
\begin{tabular}{lll}
\textbf{Projection} & \textbf{Object} & \textbf{Explicit form} \\
\hline
Closed & $\PV(Y)$ & Schouten--Nijenhuis bracket \\
Open & $D^b\mathrm{Coh}(Y)$ & Derived category \\
$(2,0)$-mixed & $R_{\mathrm{KS}}(z)$
& $\Omega_{\PV}^{\sn}/z$ (Gaussian: leading term only) \\
$(1,1)$-mixed & $\beta_{\der}$
& HKR quasi-isomorphism \\
Self-loop & $\Theta_{\mathrm{KS}}$
& $\kappa_{\mathrm{KS}}=h^{1,1}-h^{2,1}$, shadow depth~$2$
\end{tabular}
\end{center}
The convolution algebra is simpler than the $M2$ case because
the boundary algebra is commutative (polyvector fields with
wedge product) and the shadow obstruction tower terminates at degree~$2$
(Gaussian class). The MC element $\alpha_{\mathrm{KS}}$ is
therefore determined by its binary component.
\end{computation}


\subsection{Concluding identification: the correspondence as Swiss-cheese}
% label removed: subsec:sc-conclusion

The four sections of this extension have developed a single theme:
the bulk-boundary-line Koszul triangle is not a collection of
pairwise comparisons but the structure theorem for a single
two-coloured convolution algebra.

\begin{enumerate}[label=(\arabic*),leftmargin=1.8em]
\item The $M2$-brane computation
(\S\ref{sec:m2-full-computation})
demonstrated that the double-loop algebra of Costello's
construction admits a complete holographic modular Koszul datum
with explicit quantum corrections, infinite shadow depth, and
a two-parameter $R$-matrix.

\item The Kodaira--Spencer example
(\S\ref{sec:kodaira-spencer-example})
demonstrated that the triangle is a theorem (not merely a
conjecture) in the B-model setting: Koszul duality exchanges
polyvector fields and differential forms, line operators are
D-branes, and the shadow connection at genus~$0$ is the
Gauss--Manin connection.

\item The global Koszul triangle theorem
(\S\ref{sec:global-koszul-triangle})
gave the strongest current formulation, identified the defect
holonomy as the unique obstacle to globalization, and reduced
the standard-family case to the corresponding genus-$0$
holonomy packets after the structural hypotheses are checked.

\item The Swiss-cheese convolution algebra
(\S\ref{sec:sc-convolution-structural})
proved that the five structures of the triangle are projections
of a single MC element $\alpha_T$, established the precise
Steinberg-type identifications (S1)--(S7), and proved the mutual
determination theorem.
\end{enumerate}

The unifying object is the Swiss-cheese convolution $\Linf$-algebra
$\gSC_T$ and its canonical MC element $\alpha_T$. The triangle
is a shadow of the correspondence. The correspondence is the
bar complex on $\FM_k(\C)\times\Conf_k^{<}(\R)$. The presented
algebra is $A^!$. The center is the bulk. The monodromy is the
$R$-matrix. The curvature is the genus tower. All of this from
one element.

\section{Wave 13 DNA perspective}
\label{sec:ht-bbl-ext-wave13-DNA}

\begin{remark}[HT bulk--boundary--line extension at the Wave~13 base]
\label{rem:ht-bbl-ext-DNA-1}
The HT bulk--boundary--line extension programme, when specialised to
the Wave~13 K3 bialgebra $\mathbf H_{\Delta_5}$, exhibits the
three-layer structure
\begin{itemize}
\item \emph{Bulk} $=$ Hochschild cochain complex
      $C^\bullet_{\mathrm{ch}}(\mathbf H_{\Delta_5}, \mathbf H_{\Delta_5})$
      in $\mathrm{CoHA}_{K3\times E}$;
\item \emph{Boundary} $=$ $\mathbf H_{\Delta_5}$ itself, as the
      left-module category $\mathbf H_{\Delta_5}\text{-}\mathrm{mod}^{\mathrm{ch}}$
      on a boundary component of the bi-based Ran-space;
\item \emph{Line} $=$ the $24$-fold line-operator algebra on
      $\Sigma_{0,24}$, supplied by the class-$\mathcal S$ $A_1$
      construction with flavour $\mathfrak{su}(2)^{24}$.
\end{itemize}
The three layers are related by the Wave~13 bulk--boundary--line
Koszul-duality tower: Bulk $=$ centre of Boundary, Boundary $=$ centre
of Line. This is the Wave~13 realisation of the ${\sf BBL}$ operadic
correspondence developed in this chapter, with the specific $24$-fold
multiplicity forced by the Steiner $S(5,8,24)$ rigidity and the
Mukai-doubling $K^{\kappa_{\mathrm{ch}}} = 8$. Cross-ref Vol.~III,
Chapter~\ref{ch:k3-chiral-bialgebra-platonic}.
\end{remark}
